% ============================================
% Главный файл отчёта НИР
% Проект: sdCloud
% ============================================

% ============================================
% Преамбула отчёта НИР
% Проект: sdCloud — AI-ассистент для генерации
% системно-динамических моделей
% ============================================

\documentclass[14pt,a4paper]{extreport}

% ---------- Кодировка и шрифты ----------
\usepackage[utf8]{inputenc}
\usepackage[T2A]{fontenc}

% ---------- Язык ----------
\usepackage[russian]{babel}

% ---------- Геометрия страницы (ГОСТ) ----------
\usepackage{geometry}
\geometry{
    a4paper,
    left=30mm,
    right=15mm,
    top=20mm,
    bottom=20mm,
    footskip=15mm,
    headsep=10mm
}

% ---------- Интервалы ----------
\usepackage{setspace}
\onehalfspacing

% ---------- Абзацный отступ ----------
\usepackage{indentfirst}
\setlength{\parindent}{1.25cm}

% ---------- Математика ----------
\usepackage{amsmath,amssymb,amsthm}

% ---------- Графика ----------
\usepackage{graphicx}
\usepackage{tikz}
\usetikzlibrary{shapes.geometric, arrows.meta, positioning, calc, fit, backgrounds, decorations.pathreplacing}

% ---------- Таблицы ----------
\usepackage{booktabs,longtable,multirow,array,colortbl}
\usepackage{tabularx}
\newcolumntype{L}[1]{>{\raggedright\let\newline\\\arraybackslash\hspace{0pt}}m{#1}}
\newcolumntype{C}[1]{>{\centering\let\newline\\\arraybackslash\hspace{0pt}}m{#1}}
\newcolumntype{R}[1]{>{\raggedleft\let\newline\\\arraybackslash\hspace{0pt}}m{#1}}

% ---------- Листинги кода ----------
\usepackage{listings}
\usepackage{xcolor}

% ---------- Гиперссылки ----------
\usepackage[
    colorlinks=true,
    linkcolor=black,
    citecolor=black,
    urlcolor=blue,
    bookmarks=true,
    bookmarksopen=true,
    pdftitle={sdCloud: AI-ассистент для генерации системно-динамических моделей},
    pdfauthor={ИТМО}
]{hyperref}

% ---------- Колонтитулы ----------
\usepackage{fancyhdr}
\pagestyle{fancy}
\fancyhf{}
\fancyfoot[C]{\thepage}
\renewcommand{\headrulewidth}{0pt}
\renewcommand{\footrulewidth}{0pt}
\fancypagestyle{plain}{
    \fancyhf{}
    \fancyfoot[C]{\thepage}
    \renewcommand{\headrulewidth}{0pt}
    \renewcommand{\footrulewidth}{0pt}
}

% ---------- Заголовки ----------
\usepackage{titlesec}
\titleformat{\chapter}[display]
    {\normalfont\Large\bfseries\centering}
    {\chaptertitlename\ \thechapter}{14pt}{\Large}
\titlespacing*{\chapter}{0pt}{-20pt}{20pt}

\titleformat{\section}
    {\normalfont\large\bfseries}
    {\thesection}{1em}{}
\titleformat{\subsection}
    {\normalfont\normalsize\bfseries}
    {\thesubsection}{1em}{}
\titleformat{\subsubsection}
    {\normalfont\normalsize\itshape}
    {\thesubsubsection}{1em}{}

% ---------- Нумерация ----------
\renewcommand{\thechapter}{\arabic{chapter}}
\renewcommand{\thesection}{\thechapter.\arabic{section}}
\renewcommand{\thesubsection}{\thesection.\arabic{subsection}}
\renewcommand{\thesubsubsection}{\thesubsection.\arabic{subsubsection}}

% ---------- Названия элементов ----------
\addto\captionsrussian{
    \renewcommand{\contentsname}{СОДЕРЖАНИЕ}
    \renewcommand{\chaptername}{ГЛАВА}
    \renewcommand{\figurename}{Рис.}
    \renewcommand{\tablename}{Таблица}
    \renewcommand{\bibname}{СПИСОК ЛИТЕРАТУРЫ}
    \renewcommand{\lstlistingname}{Листинг}
}

% ---------- Подписи к рисункам и таблицам ----------
\usepackage{caption}
\captionsetup[figure]{labelsep=endash, justification=centering, font=small}
\captionsetup[table]{labelsep=endash, justification=justified, font=small, singlelinecheck=false}
\captionsetup[lstlisting]{labelsep=endash, font=small}

% ---------- Стили листингов ----------
\lstset{
    basicstyle=\small\ttfamily,
    keywordstyle=\color{blue}\bfseries,
    commentstyle=\color{gray}\itshape,
    stringstyle=\color{red},
    numbers=left,
    numberstyle=\tiny\color{gray},
    numbersep=8pt,
    frame=single,
    breaklines=true,
    breakatwhitespace=true,
    tabsize=4,
    showstringspaces=false,
    inputencoding=utf8,
    extendedchars=true,
    literate=
        {а}{\cyra}1 {б}{\cyrb}1 {в}{\cyrv}1 {г}{\cyrg}1 {д}{\cyrd}1
        {е}{\cyre}1 {ё}{\cyryo}1 {ж}{\cyrzh}1 {з}{\cyrz}1 {и}{\cyri}1
        {й}{\cyrishrt}1 {к}{\cyrk}1 {л}{\cyrl}1 {м}{\cyrm}1 {н}{\cyrn}1
        {о}{\cyro}1 {п}{\cyrp}1 {р}{\cyrr}1 {с}{\cyrs}1 {т}{\cyrt}1
        {у}{\cyru}1 {ф}{\cyrf}1 {х}{\cyrh}1 {ц}{\cyrc}1 {ч}{\cyrch}1
        {ш}{\cyrsh}1 {щ}{\cyrshch}1 {ъ}{\cyrhrdsn}1 {ы}{\cyry}1 {ь}{\cyrsoftsign}1
        {э}{\cyrie}1 {ю}{\cyryu}1 {я}{\cyrya}1
        {А}{\CYRA}1 {Б}{\CYRB}1 {В}{\CYRV}1 {Г}{\CYRG}1 {Д}{\CYRD}1
        {Е}{\CYRE}1 {Ё}{\CYRYO}1 {Ж}{\CYRZH}1 {З}{\CYRZ}1 {И}{\CYRI}1
        {Й}{\CYRISHRT}1 {К}{\CYRK}1 {Л}{\CYRL}1 {М}{\CYRM}1 {Н}{\CYRN}1
        {О}{\CYRO}1 {П}{\CYRP}1 {Р}{\CYRR}1 {С}{\CYRS}1 {Т}{\CYRT}1
        {У}{\CYRU}1 {Ф}{\CYRF}1 {Х}{\CYRH}1 {Ц}{\CYRC}1 {Ч}{\CYRCH}1
        {Ш}{\CYRSH}1 {Щ}{\CYRSHCH}1 {Ъ}{\CYRHRDSN}1 {Ы}{\CYRY}1 {Ь}{\CYRSOFTSIGN}1
        {Э}{\CYRIE}1 {Ю}{\CYRYU}1 {Я}{\CYRYA}1
}

% ---------- Окружения ----------
\newtheorem{definition}{Определение}[chapter]
\newtheorem{theorem}{Теорема}[chapter]

% ---------- Пакеты для алгоритмов ----------
\usepackage{algorithm}
\usepackage{algpseudocode}
\floatname{algorithm}{Алгоритм}

% ---------- Дополнительные пакеты ----------
\usepackage{enumitem}
\setlist[itemize]{leftmargin=1.5cm,itemsep=2pt}
\setlist[enumerate]{leftmargin=1.5cm,itemsep=2pt}
\usepackage{float}
\usepackage{csquotes}
\usepackage{url}
\usepackage{breakurl}
\usepackage{footnote}
\usepackage{microtype}
\usepackage{etoolbox}

% ---------- Счётчики ----------
\setcounter{secnumdepth}{3}
\setcounter{tocdepth}{2}

% ---------- Компактные списки ----------
\let\olditemize\itemize
\let\oldenumerate\enumerate
\renewenvironment{itemize}{
    \olditemize[topsep=4pt,parsep=2pt]
}{\endolditemize}
\renewenvironment{enumerate}{
    \oldenumerate[topsep=4pt,parsep=2pt]
}{\endoldenumerate}


\begin{document}

% ---------- Титульная страница ----------
% ============================================
% Титульная страница
% ============================================

\begin{titlepage}
    \thispagestyle{empty}
    \setlength{\parindent}{0pt}

    \begin{center}
        \small
        МИНИСТЕРСТВО НАУКИ И ВЫСШЕГО ОБРАЗОВАНИЯ\\
        РОССИЙСКОЙ ФЕДЕРАЦИИ\\[6pt]
        Федеральное государственное автономное образовательное\\
        учреждение высшего образования\\
        \textbf{«НАЦИОНАЛЬНЫЙ ИССЛЕДОВАТЕЛЬСКИЙ УНИВЕРСИТЕТ ИТМО»}\\
        (ФГАОУ ВО «НИУ ИТМО")
    \end{center}

    \vspace{1.5cm}

    \begin{center}
        \small
        Факультет программной инженерии и компьютерной техники\\[3pt]
        Направление подготовки 09.04.04 --- Программная инженерия
    \end{center}

    \vspace{2cm}

    \begin{center}
        \large\bfseries
        ОТЧЁТ\\
        О НАУЧНО-ИССЛЕДОВАТЕЛЬСКОЙ РАБОТЕ
    \end{center}

    \vspace{1cm}

    \begin{center}
        \large\bfseries
        по теме:\\[0.5cm]
        \Large
        МЕТОДЫ И АЛГОРИТМЫ АВТОМАТИЗИРОВАННОЙ\\
        ГЕНЕРАЦИИ СИСТЕМНО-ДИНАМИЧЕСКИХ МОДЕЛЕЙ\\
        НА ОСНОВЕ БОЛЬШИХ ЯЗЫКОВЫХ МОДЕЛЕЙ
    \end{center}

    \vspace{2cm}

    \begin{flushright}
        \begin{minipage}{0.5\textwidth}
            \small
            \begin{tabular}{@{}ll@{}}
                Выполнил: студент группы P4122 & \\[6pt]
                \underline{\hspace{4cm}} & \\[12pt]
                Руководитель: & \\[3pt]
                ст. преп., к.т.н. & \\[3pt]
                И.\,А.~Перл & \\[6pt]
                \underline{\hspace{4cm}} & \\[6pt]
                «\underline{\hspace{1cm}}» \underline{\hspace{3cm}} 2025 г.
            \end{tabular}
        \end{minipage}
    \end{flushright}

    \vfill

    \begin{center}
        \small
        Санкт-Петербург --- 2025
    \end{center}

\end{titlepage}


% ---------- Аннотация ----------
\newpage
\thispagestyle{empty}
\chapter*{АННОТАЦИЯ}
\addcontentsline{toc}{chapter}{АННОТАЦИЯ}

Отчёт содержит \pageref{LastPage}~страниц, \totalfigures~рисунков, \totaltables~таблиц, \thetotallstlistings~листингов, 32~источника.

\textbf{Ключевые слова:} системная динамика, большие языковые модели, автоматизированное моделирование, FastAPI, Next.js, генерация моделей, диаграммы причинно-следственных связей, детерминированная генерация кода, SciPy, PySD, гостевой режим.

\textbf{Объект исследования:} процессы создания системно-динамических моделей с применением методов искусственного интеллекта.

\textbf{Цель работы:} разработка методов и алгоритмов автоматизированной генерации системно-динамических моделей на основе текстовых описаний с использованием больших языковых моделей.

В работе проведён обзор предметной области, включающий теоретические основы системной динамики, программные платформы моделирования, большие языковые модели и существующие работы по их интеграции. Спроектирована и реализована программная платформа \texttt{sdCloud}, включающая backend на FastAPI, frontend на Next.js, конвейер генерации моделей через LLM, модуль валидации графовых структур, детерминированный генератор Python-кода и fallback-симуляцию через SciPy. Реализован гостевой режим с сохранением истории в браузере для демонстрации возможностей без регистрации.

\textbf{Научная новизна:} предложен многоэтапный конвейер генерации системно-динамических моделей с промежуточными представлениями CLD и SFD; разработан метод валидации графовых структур на основе топологического анализа; предложен детерминированный подход к генерации исполняемого Python-кода по структурированным уравнениям, снижающий влияние галлюцинаций LLM; создана программная платформа, интегрирующая LLM с полным циклом системно-динамического моделирования.

\textbf{Практическая значимость:} снижение барьера входа в методологию системной динамики, ускорение процесса моделирования, образовательный потенциал.

\vspace{1cm}

% ---------- Содержание ----------
\newpage
\tableofcontents

% ---------- Введение ----------
\newpage
% ============================================
% Введение
% ============================================

\chapter*{ВВЕДЕНИЕ}
\addcontentsline{toc}{chapter}{ВВЕДЕНИЕ}

\section*{Актуальность исследования}
\addcontentsline{toc}{section}{Актуальность исследования}

Системная динамика как научная дисциплина и методология моделирования сложных систем имеет глубокие исторические корни, уходящие к работам Джей Форрестера в 1960-х годах в Массачусетском технологическом институте. За прошедшие десятилетия методология зарекомендовала себя как мощный инструмент анализа нелинейного поведения сложных систем во времени, получив широкое распространение в областях управления предприятиями, экологического моделирования, анализа социально-экономических процессов и формирования государственной политики.

Традиционно создание системно-динамических моделей требует от исследователя существенных затрат времени и интеллектуальных ресурсов. Процесс включает в себя этапы качественного концептуального моделирования, построения диаграмм причинно-следственных связей (Causal Loop Diagrams, CLD), разработки диаграмм потоков и уровней (Stock and Flow Diagrams, SFD), формализации уравнений модели, параметризации, верификации и валидации. Каждый из этих этапов требует глубокого понимания как предметной области моделируемой системы, так и специфики самой методологии системной динамики.

Существенным барьером для широкого внедрения системной динамики остаётся необходимость в специализированной подготовке. Исследователям необходимо освоить специфические инструменты моделирования, такие как Vensim, Stella или AnyLogic, понять семантику конструкций системной динамики, а также развить навыки перехода от текстового описания системы к формальной модели. Это ограничивает применение методологии кругом профессиональных системных аналитиков и замедляет распространение системно-динамического подхода в смежных дисциплинах.

Параллельно с развитием методологии системной динамики в последние годы наблюдается бурный прогресс в области больших языковых моделей (Large Language Models, LLM). Модели семейства GPT, Claude, а также открытые альтернативы (LLaMA, Mistral, DeepSeek) продемонстрировали способность к пониманию сложных текстовых описаний, генерации структурированных данных и формальному представлению знаний. Эти способности создают предпосылки для использования LLM в качестве интеллектуального посредника между человеком и инструментами системной динамики.

Актуальность настоящего исследования обусловлена конвергенцией двух указанных трендов: зрелости методологии системной динамики, нуждающейся в инструментах снижения барьера входа, и достижением большими языковыми моделями уровня компетентности, достаточного для участия в процессе моделирования. Создание программной платформы, интегрирующей LLM с конвейером системно-динамического моделирования, представляет собой направление, имеющее как научную новизну, так и значительный практический потенциал.

\section*{Цель и задачи исследования}
\addcontentsline{toc}{section}{Цель и задачи исследования}

\textbf{Целью} настоящей научно-исследовательской работы является разработка методов и алгоритмов автоматизированной генерации системно-динамических моделей на основе текстовых описаний с использованием больших языковых моделей, а также создание программной платформы \texttt{sdCloud}, реализующей предложенные подходы.

Для достижения поставленной цели сформулированы следующие \textbf{задачи}:

\begin{enumerate}
    \item Провести обзор предметной области: изучить теоретические основы системной динамики, проанализировать существующие программные платформы моделирования, исследовать возможности больших языковых моделей в контексте генерации структурированных данных, а также проанализировать существующие работы по интеграции LLM и системной динамики.

    \item Спроектировать архитектуру программной платформы \texttt{sdCloud}, включающей подсистемы взаимодействия с пользователем, управления сессиями моделирования, интеграции с LLM, генерации и валидации моделей, а также хранения данных.

    \item Реализовать backend-подсистему на базе фреймворка FastAPI с использованием реляционной базы данных PostgreSQL, кэширования Redis и асинхронной очереди задач Taskiq.

    \item Реализовать frontend-приложение на базе фреймворка Next.js с применением методологии Feature-Sliced Design (FSD) и библиотеки компонентов shadcn/ui.

    \item Разработать конвейер генерации системно-динамических моделей, включающий этапы извлечения переменных, создания концептуальной карты (CLD), формирования системы уравнений, построения диаграммы потоков и уровней (SFD), а также генерации исполняемого Python-кода и Vensim-экспорта.

    \item Реализовать модуль валидации графовых структур моделей с использованием библиотеки NetworkX для проверки корректности топологии сгенерированных диаграмм и согласованности переменных, уравнений и CLD.

    \item Разработать детерминированный генератор Python-кода по структурированным уравнениям и модуль fallback-симуляции на базе SciPy \texttt{solve\_ivp} для случаев, когда прямая загрузка в PySD невозможна.

    \item Реализовать гостевой режим работы платформы, позволяющий анонимным пользователям генерировать модели и сохранять историю взаимодействия в браузере без обращений к серверной базе данных.

    \item Провести тестирование и оценку работоспособности разработанной платформы на типовых задачах системной динамики.
\end{enumerate}

\section*{Объект и предмет исследования}
\addcontentsline{toc}{section}{Объект и предмет исследования}

\textbf{Объектом исследования} являются процессы создания и модификации системно-динамических моделей с применением методов искусственного интеллекта.

\textbf{Предметом исследования} выступают методы и алгоритмы автоматизированной генерации системно-динамических моделей на основе текстовых описаний с использованием больших языковых моделей, а также архитектурные решения программных платформ, поддерживающих подобную генерацию.

\section*{Научная новизна}
\addcontentsline{toc}{section}{Научная новизна}

Научная новизна работы заключается в следующем:

\begin{itemize}
    \item Предложен многоэтапный конвейер генерации системно-динамических моделей из естественно-языковых описаний, включающий промежуточные представления в виде CLD и SFD, что обеспечивает прозрачность и интерпретируемость процесса генерации.

    \item Разработан метод валидации графовых структур системно-динамических моделей, основанный на анализе топологических свойств графа причинно-следственных связей с использованием алгоритмов на графах, а также на проверке согласованности между переменными, уравнениями и CLD.

    \item Предложен детерминированный метод генерации исполняемого Python-кода по структурированным уравнениям модели, заменяющий нестабильную генерацию кода непосредственно LLM и снижающий влияние галлюцинаций языковой модели на корректность симуляции.

    \item Разработан fallback-механизм симуляции на основе SciPy \texttt{solve\_ivp}, позволяющий выполнять модели в случаях несовместимости сгенерированного кода с текущей версией PySD.

    \item Создана программная платформа \texttt{sdCloud}, интегрирующая современные веб-технологии, базы данных, асинхронную обработку задач, API больших языковых моделей и гостевой режим работы в единое решение для AI-ассистированного системно-динамического моделирования.
\end{itemize}

\section*{Практическая значимость}
\addcontentsline{toc}{section}{Практическая значимость}

Практическая значимость результатов работы определяется следующими факторами:

\begin{itemize}
    \item Разработанная платформа \texttt{sdCloud} может быть использована в образовательном процессе для обучения основам системной динамики, позволяя студентам фокусироваться на анализе моделей, а не на технических аспектах их создания.

    \item Платформа снижает барьер входа в методологию системной динамики для специалистов, не имеющих специализированной подготовки в области моделирования, в том числе благодаря гостевому режиму без обязательной регистрации.

    \item Детерминированный генератор Python-кода и fallback-симуляция через SciPy повышают надёжность и воспроизводимость результатов по сравнению с прямой генерацией исполняемого кода LLM.

    \item Исходный код платформы открыт и может служить основой для дальнейших исследований в области автоматизированного моделирования сложных систем.

    \item Архитектурные решения, применённые в проекте, могут быть адаптированы для создания аналогичных систем в смежных предметных областях.
\end{itemize}

\section*{Структура отчёта}
\addcontentsline{toc}{section}{Структура отчёта}

Отчёт о научно-исследовательской работе состоит из введения, четырёх глав, заключения и списка литературы.

В \textbf{первой главе} представлен обзор литературы по теоретическим основам системной динамики, существующим программным платформам моделирования, большим языковым моделям и существующим работам по интеграции LLM и системной динамики.

Во \textbf{второй главе} описана архитектура программной платформы \texttt{sdCloud}, включая подсистемы backend, frontend, модули генерации и валидации моделей, интеграцию с внешними сервисами LLM, гостевой режим и механизм fallback-симуляции.

В \textbf{третьей главе} приведены детали реализации платформы, включая структуру проекта, описание API, моделей данных, детерминированного генератора Python-кода, SciPy-симуляции, постпроцессинга уравнений, процесса генерации и результаты тестирования.

В \textbf{четвёртой главе} обсуждаются ограничения текущей реализации (включая ограничения бесплатных языковых моделей), планируемые направления развития и потенциальные области применения платфорформы.

В \textbf{заключении} подводятся итоги работы и формулируются выводы.


% ---------- Глава 1 ----------
\newpage
% ============================================
% Глава 1. Обзор литературы
% ============================================

\chapter{ОБЗОР ПРЕДМЕТНОЙ ОБЛАСТИ}

\section{Системная динамика: методология и приложения}

\subsection{Исторические предпосылки и фундаментальные основы}

Системная динамика представляет собой методологию моделирования и анализа поведения сложных систем, характеризующихся наличием нелинейных взаимосвязей, обратных связей и запаздываний. Методология была основана Джейом Форрестером в 1950--1960-х годах в Массачусетском технологическом институте. Первоначально разработанная для анализа промышленных систем, методология быстро нашла применение вUrban Dynamics \cite{forrester1969urban}, изучении глобальных проблем \cite{meadows1972limits} и бизнес-стратегиях \cite{sterman2000business}.

Фундаментальные принципы системной динамики были изложены в работе \cite{forrester1961industrial}, где Форрестер ввёл ключевые понятия: \emph{уровни} (stocks, накопители), \emph{потоки} (flows) и \emph{обратные связи} (feedback loops). Уровни представляют собой состояния системы, характеризующиеся наличием памяти --- они аккумулируют значения на протяжении времени. Потоки определяют скорости изменения уровней и могут зависеть от текущих значений уровней, внешних параметров и времени. Обратные связи формируют структуру взаимодействия элементов системы и определяют её поведение во времени.

Поведение системно-динамической модели определяется решением системы обыкновенных дифференциальных уравнений первого порядка, которые в общем виде записываются следующим образом:

\begin{equation}
    \frac{dS_i}{dt} = F_i(S_1, S_2, \ldots, S_n; P_1, P_2, \ldots, P_m; t), \quad i = 1, \ldots, n
    \label{eq:sde_general}
\end{equation}

где $S_i$ --- уровни (state variables) модели, $F_i$ --- потоки, $P_j$ --- внешние параметры, $t$ --- время, $n$ --- число уровней, $m$ --- число параметров.

Форрестер выделил два фундаментальных типа петель обратной связи: \emph{усиливающие} (reinforcing, положительная обратная связь) и \emph{балансирующие} (balancing, отрицательная обратная связь). Усиливающие петли порождают экспоненциальный рост или коллапс --- поведение, при котором небольшие возмущения усиливаются системой. Балансирующие петли, напротив, стремятся привести систему к равновесному состоянию, компенсируя отклонения от целевого значения. Взаимодействие этих двух типов петель определяет основные архетипы поведения системы.

\subsection{Концептуальное моделирование: диаграммы причинно-следственных связей}

Перед формализацией математической модели исследователь выполняет качественный этап концептуального моделирования, результатом которого является \emph{диаграмма причинно-следственных связей} (Causal Loop Diagram, CLD). CLD представляет собой ориентированный граф, вершины которого соответствуют переменным системы, а рёбра --- причинно-следственным влияниям между переменными.

Каждое ребро графа помечается знаком: плюс (+) означает, что увеличение причины ведёт к увеличению следствия (при прочих равных), минус (-) означает, что увеличение причины ведёт к уменьшению следствия. Петля обратной связи классифицируется как усиливающая, если число отрицательных связей в ней чётно, и как балансирующая, если число отрицательных связей нечётно.

CLD служит важнейшим инструментом коммуникации между участниками моделирования, поскольку позволяет представить структуру системы в наглядной форме, доступной для понимания специалистами без математической подготовки. Кроме того, CLD фиксирует гипотезы исследователя о структуре системы и может использоваться для структурного анализа до перехода к количественному моделированию.

\subsection{Диаграммы потоков и уровней}

На следующем этапе моделирования CLD трансформируется в \emph{диаграмму потоков и уровней} (Stock and Flow Diagram, SFD). SFD является более формальным представлением модели, непосредственно отражающим структуру системы дифференциальных уравнений. В SFD вводятся специальные графические элементы:

\begin{itemize}
    \item \textbf{Уровни} (stocks) --- изображаются в виде прямоугольников и представляют накопители, интегрирующие потоки;
    \item \textbf{Потоки} (flows) --- изображаются в виде клапанов (стрелок с регулятором) и определяют скорость изменения уровней;
    \item \textbf{Вспомогательные переменные} (auxiliaries) --- изображаются в виде окружностей и представляют промежуточные вычисления;
    \item \textbf{Константы и экзогенные переменные} --- также изображаются в виде окружностей или ромбов;
    \item \textbf{Источники и стоки} (sources/sinks) --- изображаются в виде облаков и представляют границы моделируемой системы.
\end{itemize}

Каждый уровень в SFD соответствует переменной состояния в системе уравнений, а каждый поток --- правой части дифференциального уравнения. Топология SFD однозначно определяет структуру системы уравнений, хотя конкретные функциональные зависимости должны быть заданы отдельно.

\subsection{Применения системной динамики}

Системная динамика нашла применение в широком спектре предметных областей. Стерман \cite{sterman2000business} систематизировал применение методологии в бизнес-стратегиях, операционном управлении и управлении цепями поставок. В области экологии работы \cite{meadows1972limits} продемонстрировали способность системно-динамических моделей описывать глобальные процессы и выявлять режимы устойчивого развития.

В сфере здравоохранения системная динамика используется для моделирования распространения инфекционных заболеваний, планирования ресурсов медицинских учреждений и оценки вмешательств в системы общественного здравоохранения \cite{rahimi2020healthcare}. В энергетике методология применяется для анализа перехода к возобновляемым источникам энергии и моделирования энергетических рынков \cite{keith2017energy}.

Несмотря на широкий спектр приложений, методология сталкивается с ограничениями, связанными с трудоёмкостью процесса моделирования, необходимостью в специализированной экспертизе и сложностью валидации моделей. Именно эти ограничения определяют потенциальную ценность автоматизации с использованием современных методов искусственного интеллекта.


\section{Программные платформы системной динамики}

\subsection{Коммерческие системы моделирования}

Современный ландшафт инструментов системной динамики включает как коммерческие, так и открытые программные продукты. Наиболее широкое распространение получили следующие системы.

\textbf{Vensim} (Ventana Systems) является одним из наиболее авторитетных инструментов в сообществе системной динамики. Система поддерживает все этапы моделирования --- от построения CLD до симуляции и анализа чувствительности. Vensim реализует собственный язык моделирования и формат хранения моделей (файлы с расширением \texttt{.mdl}), который фактически является стандартом де-факто в области. Расширенные версии (Vensim DSS и Professional) включают инструменты калибровки моделей по данным, оптимизации и анализа чувствительности.

\textbf{Stella} (isee systems) --- один из старейших инструментов системной динамики, изначально разработанный для платформы Macintosh. Stella предлагает визуальный интерфейс моделирования с богатыми возможностями построения диаграмм и интерактивного анализа. Система позиционируется как образовательный инструмент и широко используется в академических учреждениях.

\textbf{AnyLogic} (The AnyLogic Company) представляет собой мультиметодную платформу моделирования, поддерживающую системную динамику наряду с агентным и дискретно-событийным моделированием. AnyLogic обеспечивает интеграцию различных парадигм моделирования в рамках единого проекта, что позволяет создавать гибридные модели сложных систем.

\textbf{Powersim Studio} --- коммерческая платформа, ориентированная на построение моделей для бизнес-анализа и принятия решений. Система включает инструменты разработки интерактивных панелей управления (dashboards) и интеграции с внешними источниками данных.

\subsection{Открытые инструменты и библиотеки}

Важным направлением развития инструментария системной динамики является создание открытых программных решений, обеспечивающих доступность методологии и возможность интеграции с другими системами.

\textbf{PySD} \cite{pysd2021} --- библиотека на языке Python, разработанная Хинкельманом и др., обеспечивающая парсинг, трансляцию и симуляцию моделей системной динамики, созданных в Vensim и Stella. PySD транслирует модели из формата \texttt{.mdl} в код Python, используя библиотеку \texttt{pandas} для представления временных рядов и \texttt{scipy} для численного интегрирования. Ключевые возможности PySD включают:

\begin{itemize}
    \item загрузку моделей Vensim (формат \texttt{.mdl}) и XMILE (формат \texttt{.stmx});
    \item симуляцию моделей с использованием различных методов интегрирования (Эйлера, Рунге-Кутта);
    \item выполнение параметрического анализа и анализа чувствительности;
    \item экспорт моделей в формате Python для интеграции с экосистемой научных вычислений.
\end{itemize}

PySD является ключевым компонентом программной инфраструктуры проекта \texttt{sdCloud}, поскольку обеспечивает возможность симуляции сгенерированных LLM моделей без необходимости использования коммерческого ПО.

\textbf{BPTK-Py} --- фреймворк на Python для системной динамики и агентного моделирования, разработанный компанией transentis. Фреймворк поддерживает моделирование в духе System Dynamics с использованием современных возможностей Python.

\textbf{Mesa} --- библиотека для агентного моделирования на Python, которая может использоваться в сочетании с компонентами системной динамики для создания гибридных моделей.

\subsection{Форматы представления моделей}

Важным аспектом интеграции различных инструментов системной динамики является унификация форматов представления моделей. Формат \textbf{XMILE} (XML Model Interchange Language) \cite{xmile2015} был разработан как стандарт обмена моделями системной динамики на базе XML. XMILE позволяет переносить модели между различными программными продуктами и обеспечивает машинную читаемость структуры модели.

Формат \textbf{MDL} (Vensim Model) использует текстовое представление, где уравнения модели записываются в декларативном стиле. Каждое уравнение имеет вид:

\begin{lstlisting}[language={}, caption={Пример записи уравнения в формате MDL}]
Variable = Expression
~ Units | Commentary
\end{lstlisting}

где \texttt{Expression} может содержать ссылки на другие переменные, функции (INTEG, SMOOTH, DELAY и др.) и константы.

Наличие стандартизированных форматов представления моделей является важным фактором для автоматизированной генерации, поскольку позволяет формализовать целевой формат выходных данных LLM.


\section{Большие языковые модели}

\subsection{Эволюция языковых моделей}

Большие языковые модели (LLM) представляют собой класс нейросетевых архитектур, основанных на механизме внимания (attention mechanism) \cite{vaswani2017attention}, обученных на массивных корпусах текстовых данных для моделирования последовательностей естественного языка. Ключевым событием в развитии LLM стала публикация архитектуры GPT (Generative Pre-trained Transformer) \cite{radford2018gpt} и её последующие итерации: GPT-2 \cite{radford2019gpt2}, GPT-3 \cite{brown2020gpt3} и GPT-4 \cite{openai2023gpt4}.

Архитектура трансформера основана на механизме самовнимания (self-attention), который позволяет модели устанавливать зависимости между элементами входной последовательности независимо от их расстояния. Формально, для входной последовательности $\mathbf{X} = (\mathbf{x}_1, \ldots, \mathbf{x}_n)$ вычисляются запросы (queries), ключи (keys) и значения (values):

\begin{equation}
    \mathbf{Q} = \mathbf{X}\mathbf{W}_Q, \quad \mathbf{K} = \mathbf{X}\mathbf{W}_K, \quad \mathbf{V} = \mathbf{X}\mathbf{W}_V
\end{equation}

где $\mathbf{W}_Q$, $\mathbf{W}_K$, $\mathbf{W}_V$ --- обучаемые матрицы проекции. Веса внимания вычисляются как:

\begin{equation}
    \text{Attention}(\mathbf{Q}, \mathbf{K}, \mathbf{V}) = \text{softmax}\left(\frac{\mathbf{Q}\mathbf{K}^T}{\sqrt{d_k}}\right)\mathbf{V}
\end{equation}

где $d_k$ --- размерность пространства ключей, а нормировка на $\sqrt{d_k}$ предотвращает чрезмерно малые градиенты при больших размерностях.

Масштабирование моделей --- увеличение числа параметров, объёма обучающих данных и вычислительных ресурсов --- привело к качественному скачку в способностях LLM. Модели семейства GPT-4, Claude 3 (Anthropic), Gemini (Google) продемонстрировали способность к выполнению сложных интеллектуальных задач, включая рассуждения, планирование, генерацию кода и решение математических задач.

\subsection{Стратегии использования LLM}

В контексте прикладных задач выделяются несколько стратегий использования LLM:

\textbf{Zero-shot prompting} предполагает подачу модели запроса без примеров выполнения задачи. Модель должна опираться исключительно на знания, полученные в процессе предварительного обучения. Этот подход прост в реализации, но часто даёт нестабильные результаты для специализированных задач.

\textbf{Few-shot prompting} \cite{brown2020gpt3} включает в запрос несколько примеров входных данных и ожидаемых ответов. Демонстрация примеров позволяет модели выявить шаблон задачи и сформировать ответ в соответствии с ожидаемым форматом. Этот подход существенно повышает качество генерации структурированных данных.

\textbf{Chain-of-Thought (CoT) prompting} \cite{wei2022chain} предполагает стимулирование модели к пошаговому рассуждению перед формированием финального ответа. В запросе явно предлагается модели «рассуждать шаг за шагом» (\textit{«think step by step»}), что приводит к существенному улучшению качества ответов на задачи, требующие многошагового логического вывода.

\textbf{ReAct (Reasoning and Acting)} \cite{yao2023react} комбинирует генерацию рассуждений с выполнением действий (вызовом инструментов), что позволяет модели взаимодействовать с внешними системами и получать дополнительную информацию в процессе решения задачи.

\textbf{Retrieval-Augmented Generation (RAG)} \cite{lewis2020retrieval} дополняет генерацию модели извлечением релевантной информации из внешних источников знаний. В контексте системной динамики RAG может использоваться для извлечения примеров моделей из базы знаний, что обогащает контекст генерации.

\subsection{Открытые языковые модели}

Наряду с проприетарными моделями активно развивается экосистема открытых LLM. Модели семейства \textbf{LLaMA} (Meta) \cite{touvron2023llama}, \textbf{Mistral} (Mistral AI) \cite{jiang2023mistral}, \textbf{Falcon} (TII), \textbf{DeepSeek} (DeepSeek AI) и другие демонстрируют качество генерации, сопоставимое с проприетарными аналогами при определённых условиях.

Открытые модели имеют ключевое преимущество для проектов, связанных с обработкой конфиденциальных данных или требующих развёртывания в закрытой инфраструктуре. В контексте проекта \texttt{sdCloud} открытые модели рассматриваются как перспективное направление для снижения зависимости от внешних API и обеспечения полного контроля над конвейером генерации.


\section{AI-ассистенты в системно-динамическом моделировании}

\subsection{Обзор существующих решений}

Интеграция больших языковых моделей с инструментами системной динамики является активно развивающимся направлением исследований. В последние годы появился ряд работ, исследующих различные аспекты этой интеграции.

Работа \cite{hu2025chatpysd} представляет систему \textbf{ChatPySD} --- плагин для ChatGPT, обеспечивающий загрузку, симуляцию и анализ моделей системной динамики непосредственно в интерфейсе чат-бота. ChatPySD использует PySD как движок симуляции и демонстрирует возможность естественно-языкового взаимодействия с существующими моделями, включая изменение параметров, запуск симуляций и визуализацию результатов. Ключевым ограничением ChatPySD является фокус на взаимодействии с готовыми моделями, а не на генерации новых моделей из текстовых описаний.

В работе \cite{akhavan2024generative} проведён обзор возможностей генеративного ИИ в контексте имитационного моделирования, включая системную динамику. Авторы выделяют три направления применения: генерация кода моделей, генерация входных данных для симуляций и генерация описаний результатов моделирования. Работа подчёркивает потенциал LLM для автоматизации рутинных задач в процессе моделирования, но отмечает необходимость человеческой верификации сгенерированных результатов.

Исследование \cite{liu2025leveraging} направлено на использование LLM для автоматической генерации CLD из текстовых описаний. Авторы демонстрируют, что GPT-4 способен извлекать переменные и причинно-следственные связи из текстовых описаний систем и формировать корректные CLD для относительно простых случаев. Качество генерации существенно зависит от сложности описываемой системы и детальности текстового описания.

\subsection{Конвейеры генерации моделей}

Совокупность существующих работ позволяет сформулировать общую архитектуру конвейера генерации системно-динамических моделей с использованием LLM. Типичный конвейер включает следующие этапы:

\begin{enumerate}
    \item \textbf{Извлечение переменных} (Variable Extraction) --- идентификация релевантных переменных из текстового описания системы.
    \item \textbf{Построение CLD} --- определение причинно-следственных связей между извлечёнными переменными и формирование петель обратной связи.
    \item \textbf{Трансляция в SFD} --- преобразование CLD в диаграмму потоков и уровней с явным выделением накопителей и потоков.
    \item \textbf{Генерация уравнений} --- формализация функциональных зависимостей в виде системы уравнений системной динамики.
    \item \textbf{Параметризация} --- назначение численных значений параметрам и начальных условий.
    \item \textbf{Верификация и валидация} --- проверка корректности сгенерированной модели и сопоставление её поведения с ожидаемым.
\end{enumerate}

Каждый этап конвейера может выполняться различными LLM с использованием специализированных промптов. Важным аспектом является итеративный характер процесса: результаты каждого этапа могут корректироваться на основе обратной связи от последующих этапов или пользователя.

\subsection{Проблемы и вызовы}

Несмотря на впечатляющие результаты отдельных исследований, область интеграции LLM и системной динамики сталкивается с рядом фундаментальных проблем:

\textbf{Галлюцинации моделей} --- LLM могут генерировать структурно корректные, но содержательно ошибочные модели. Модель может ввести переменные, не имеющие отношения к описываемой системе, или установить некорректные причинно-следственные связи.

\textbf{Сложность иерархических систем} --- для систем с большим числом переменных и сложной структурой обратных связей качество генерации существенно снижается. LLM имеют ограничения на контекст, которые могут препятствовать генерации комплексных моделей.

\textbf{Недетерминированность} --- стохастическая природа генерации LLM приводит к тому, что повторный запрос с идентичным промптом может дать различные результаты, что затрудняет воспроизводимость.

\textbf{Верификация и валидация} --- автоматическая проверка корректности сгенерированных моделей остаётся нерешённой проблемой. Требуются специализированные методы валидации графовых структур и уравнений моделей.

Проект \texttt{sdCloud} направлен на систематический подход к решению указанных проблем через внедрение многоэтапного конвейера генерации с промежуточной валидацией и использованием специализированных алгоритмов проверки топологии моделей.


% ---------- Глава 2 ----------
\newpage
% ============================================
% Глава 2. Проектирование системы sdCloud
% ============================================

\chapter{ПРОЕКТИРОВАНИЕ СИСТЕМЫ SDCLOUD}

\section{Общая архитектура платформы}

\subsection{Принципы проектирования}

Проектирование программной платформы \texttt{sdCloud} основано на совокупности архитектурных принципов, обеспечивающих масштабируемость, поддерживаемость и расширяемость системы:

\begin{itemize}
    \item \textbf{Разделение ответственности} (Separation of Concerns) --- чёткое разграничение функциональности между компонентами системы: frontend отвечает за представление и взаимодействие с пользователем, backend реализует бизнес-логику, а внешние сервисы предоставляют специализированные возможности (LLM, хранение данных).
    \item \textbf{Многоуровневая архитектура} (Layered Architecture) --- backend построен по принципу разделения на слои: слой API (контроллеры), слой сервисов (бизнес-логика), слой доступа к данным (репозитории) и слой моделей данных.
    \item \textbf{Асинхронная обработка} --- длительные операции (генерация моделей через LLM) выполняются асинхронно через очередь задач, что обеспечивает отзывчивость интерфейса и масштабируемость системы.
    \item \textbf{Модульность} --- система декомпозирована на функциональные модули с чётко определёнными интерфейсами, что упрощает тестирование и позволяет заменять реализации отдельных компонентов.
\end{itemize}

\subsection{Компонентная архитектура}

Платформа \texttt{sdCloud} состоит из следующих основных компонентов (рис.~\ref{fig:architecture}):

\begin{figure}[H]
    \centering
    \begin{tikzpicture}[
        node distance=0.8cm and 1.2cm,
        box/.style={draw, rounded corners=4pt, minimum width=3.2cm, minimum height=0.9cm, align=center, font=\small},
        layer/.style={draw, dashed, rounded corners=6pt, inner sep=8pt},
        arrow/.style={-{Stealth[scale=0.8]}, thick}
    ]
        % Client Layer
        \node[box, fill=blue!10] (client) {Веб-браузер\\(Next.js SPA)};

        % API Gateway Layer
        \node[box, fill=orange!10, below=of client] (api) {FastAPI Gateway\\REST API};

        % Service Layer
        \node[box, fill=green!10, below left=0.8cm and 0.3cm of api] (session) {Сервис\\сессий};
        \node[box, fill=green!10, below=of api] (gen) {Сервис\\генерации};
        \node[box, fill=green!10, below right=0.8cm and 0.3cm of api] (validate) {Сервис\\валидации};

        % Task Queue
        \node[box, fill=purple!10, below=1.5cm of gen] (taskiq) {Taskiq\\(очередь задач)};

        % Data Layer
        \node[box, fill=gray!15, below left=0.8cm and 0.2cm of taskiq] (postgres) {PostgreSQL\\(данные)};
        \node[box, fill=gray!15, below=of taskiq] (redis) {Redis\\(кэш, состояния)};

        % External
        \node[box, fill=red!10, right=2.5cm of api] (openrouter) {OpenRouter API\\(LLM)};
        \node[box, fill=red!10, below=of openrouter] (pysd) {PySD\\(симуляция)};

        % Arrows
        \draw[arrow] (client) -- (api);
        \draw[arrow] (api) -- (session);
        \draw[arrow] (api) -- (gen);
        \draw[arrow] (api) -- (validate);
        \draw[arrow] (gen) -- (taskiq);
        \draw[arrow] (taskiq) -- (postgres);
        \draw[arrow] (taskiq) -- (redis);
        \draw[arrow, dashed] (gen.east) -- ++(0.5,0) |- (openrouter.west);
        \draw[arrow, dashed] (taskiq.east) -- ++(0.5,0) |- (pysd.west);
        \draw[arrow] (session.south) -- ++(0,-0.4) -| (postgres.north);

    \end{tikzpicture}
    \caption{Архитектура платформы sdCloud}
    \label{fig:architecture}
\end{figure}

\begin{enumerate}
    \item \textbf{Веб-клиент} (Frontend) --- одностраничное приложение на базе Next.js, реализующее пользовательский интерфейс для создания сессий моделирования, ввода текстовых описаний, визуализации сгенерированных моделей (CLD, SFD, уравнения, графики симуляций) и управления результатами. Поддерживаются как авторизованный, так и гостевой режимы.

    \item \textbf{API Gateway} --- REST API на базе FastAPI, обеспечивающий точки входа для всех операций платформы, аутентификацию на основе JWT, валидацию входных данных и маршрутизацию запросов к внутренним сервисам.

    \item \textbf{Сервис сессий} --- управление сессиями моделирования, включающее создание, чтение, обновление и удаление сессий (CRUD-операции), а также управление сообщениями в рамках каждой сессии.

    \item \textbf{Сервис генерации} --- реализация конвейера генерации системно-динамических моделей, включающего взаимодействие с LLM через OpenRouter API и последовательное формирование переменных, CLD, уравнений, SFD, Python-кода и Vensim-экспорта.

    \item \textbf{Сервис валидации} --- проверка корректности графовых структур сгенерированных моделей с использованием алгоритмов на графах (библиотека NetworkX) и проверка согласованности между переменными, уравнениями и CLD.

    \item \textbf{Детерминированный генератор Python-кода} --- модуль, формирующий исполняемый Python-код модели по структурированным уравнениям без повторного обращения к LLM.

    \item \textbf{Симуляционный движок} --- PySD используется в качестве основного средства симуляции; при невозможности загрузки сгенерированного кода в PySD (например, из-за изменений в API версии 3.x) автоматически применяется fallback на SciPy \texttt{solve\_ivp}.

    \item \textbf{PostgreSQL} --- реляционная база данных для персистентного хранения пользователей, сессий, сообщений, моделей и результатов симуляций.

    \item \textbf{Redis} --- хранилище кэша и состояний, используемое для кэширования часто запрашиваемых данных.

    \item \textbf{OpenRouter API} --- внешний сервис-агрегатор, предоставляющий унифицированный доступ к различным LLM через единый REST API; взаимодействие реализовано через адаптер на базе \texttt{litellm} и OpenAI SDK.

    \item \textbf{Гостевой режим} --- набор stateless API-эндпоинтов, позволяющих анонимным пользователям синхронно генерировать модели и запускать симуляции без записи данных в PostgreSQL.
\end{enumerate}

\subsection{Взаимодействие компонентов}

Основной сценарий взаимодействия компонентов при генерации модели реализуется следующим образом:

\begin{enumerate}
    \item Пользователь вводит текстовое описание системы через веб-интерфейс.
    \item Frontend отправляет HTTP-запрос к API Gateway (авторизованный или гостевой).
    \item Для авторизованных пользователей API создаёт фоновую \texttt{asyncio}-задачу генерации; для гостей генерация выполняется синхронно в рамках одного запроса.
    \item Воркер генерации последовательно выполняет этапы конвейера, обращаясь к OpenRouter API для извлечения переменных, генерации CLD и уравнений, и строит SFD на основе уравнений.
    \item После генерации CLD и SFD выполняется валидация графовых структур и согласованности модели.
    \item Детерминированный генератор формирует Python-код модели; отдельный модуль экспортирует модель в формат Vensim .mdl.
    \item Сгенерированный Python-код передаётся в PySD; при ошибках загрузки автоматически используется SciPy \texttt{solve\_ivp} fallback.
    \item Результаты сохраняются в PostgreSQL (для авторизованных пользователей) или возвращаются напрямую в ответе (гостевой режим).
    \item Frontend отображает CLD, SFD, уравнения, Python/Vensim код и результаты симуляции.
\end{enumerate}


\section{Backend: реализация серверной части}

\subsection{Фреймворк FastAPI}

Серверная часть платформы реализована на базе фреймворка \textbf{FastAPI} \cite{fastapi2024} --- современного веб-фреймворка для Python, основанного на типизации (type hints) и стандарте ASGI. FastAPI обеспечивает следующие преимущества для проекта:

\begin{itemize}
    \item \textbf{Автоматическая валидация данных} на основе аннотаций типов Python и библиотеки Pydantic, что сокращает объём шаблонного кода и повышает надёжность.
    \item \textbf{Автоматическая генерация документации} API в форматах OpenAPI и Swagger UI, упрощающая разработку и тестирование.
    \item \textbf{Нативная поддержка асинхронного программирования} (\texttt{async}/\texttt{await}), что критично для эффективной работы с внешними API и базами данных.
    \item \textbf{Высокая производительность} --- FastAPI демонстрирует одни из лучших показателей пропускной способности среди Python-фреймворков.
\end{itemize}

\subsection{ORM SQLAlchemy и миграции Alembic}

Для взаимодействия с реляционной базой данных используется \textbf{SQLAlchemy} 2.0 \cite{sqlalchemy2024} --- наиболее зрелый и функциональный ORM-фреймворк для Python. SQLAlchemy 2.0 предоставляет:

\begin{itemize}
    \item Декларативное описание моделей данных с поддержкой типизации.
    \item Мощный язык запросов, генерирующий оптимальный SQL-код.
    \item Поддержку асинхронных сессий (\texttt{AsyncSession}) для неблокирующего доступа к базе данных.
    \item Встроенную систему миграций через интеграцию с Alembic.
\end{itemize}

\textbf{Alembic} используется для управления схемой базы данных. Система миграций обеспечивает версионирование структуры БД, возможность наката и отката изменений, а также согласованность схемы между различными окружениями (разработка, тестирование, продакшн).

\subsection{База данных PostgreSQL}

\textbf{PostgreSQL} \cite{postgresql2024} выбран в качестве основной системы управления базами данных благодаря следующим характеристикам:

\begin{itemize}
    \item Полная поддержка стандарта SQL и расширенные возможности (JSON/JSONB, массивы, пользовательские типы).
    \item Высокая надёжность и производительность.
    \item Поддержка асинхронных подключений через \texttt{asyncpg}.
    \item Расширяемость и богатая экосистема.
\end{itemize}

В контексте проекта \texttt{sdCloud} PostgreSQL хранит информацию о пользователях, сессиях моделирования, сообщениях, сгенерированных моделях (CLD, SFD, уравнения), результатах симуляций и метаданных.

\subsection{Redis}

\textbf{Redis} \cite{redis2024} используется в двух качествах:

\begin{enumerate}
    \item \textbf{Хранилище состояний задач} --- отслеживание статуса асинхронных задач генерации (в очереди, в процессе, завершено, ошибка).
    \item \textbf{Кэш} --- кэширование часто запрашиваемых данных (например, списка сессий пользователя) для снижения нагрузки на PostgreSQL.
\end{enumerate}

Асинхронный клиент \texttt{redis-py} обеспечивает неблокирующее взаимодействие с хранилищем в контексте async-приложения FastAPI.

\subsection{Асинхронная очередь задач Taskiq}

\textbf{Taskiq} \cite{taskiq2024} --- современная асинхронная очередь задач для Python, выбранная в качестве альтернативы Celery. Ключевые преимущества Taskiq:

\begin{itemize}
    \item Нативная поддержка \texttt{async}/\texttt{await} без необходимости в синхронных обёртках.
    \item Интеграция с FastAPI через middleware.
    \item Гибкая система брокеров (Redis, RabbitMQ и др.).
    \item Поддержка периодических задач (cron).
    \item Типизированные задачи с автоматической сериализацией.
\end{itemize}

В проекте \texttt{sdCloud} Taskiq используется для асинхронного выполнения конвейера генерации моделей. Каждый запрос на генерацию создаёт задачу, которая выполняется воркером независимо от HTTP-запроса пользователя.


\section{Frontend: клиентская часть}

\subsection{Фреймворк Next.js}

Клиентская часть реализована на базе \textbf{Next.js} \cite{nextjs2024} --- React-фреймворка, предоставляющего серверный рендеринг (SSR), статическую генерацию (SSG) и клиентскую навигацию. В контексте проекта \texttt{sdCloud} Next.js используется в режиме одностраничного приложения (SPA) с клиентской маршрутизацией и динамической загрузкой данных.

Выбор Next.js обусловлен следующими факторами:

\begin{itemize}
    \item Развитая экосистема инструментов и библиотек для React.
    \item Оптимизации производительности (код-сплиттинг, предзагрузка страниц).
    \item Гибкость в выборе стратегии рендеринга.
    \item Поддержка API Routes для серверных функций.
\end{itemize}

\subsection{Методология Feature-Sliced Design}

Архитектура frontend-приложения построена в соответствии с методологией \textbf{Feature-Sliced Design} (FSD) \cite{fsd2024} --- современным подходом к организации кода frontend-приложений, ориентированным на разделение по функциональным возможностям (фичам).

FSD определяет иерархическую структуру слоёв приложения:

\begin{table}[H]
    \centering
    \caption{Слои архитектуры FSD в проекте sdCloud}
    \label{tab:fsd-layers}
    \begin{tabular}{|p{3cm}|p{5cm}|p{5.5cm}|}
        \hline
        \textbf{Слой} & \textbf{Назначение} & \textbf{Примеры компонентов} \\
        \hline
        app & Инициализация приложения, провайдеры, глобальные стили & \texttt{providers/, router/, store/} \\
        \hline
        pages & Страницы приложения, композиция виджетов & \texttt{HomePage/, SessionPage/, ModelsPage/} \\
        \hline
        widgets & Самостоятельные блоки интерфейса & \texttt{ChatWidget/, ModelViewer/, Sidebar/} \\
        \hline
        features & Функциональные модули, реализующие бизнес-логику & \texttt{auth/, modelGenerate/, sessionCreate/} \\
        \hline
        entities & Бизнес-сущности предметной области & \texttt{session/, model/, user/} \\
        \hline
        shared & Переиспользуемые компоненты, утилиты, конфигурация & \texttt{api/, ui/, lib/, config/} \\
        \hline
    \end{tabular}
\end{table}

Ключевым правилом FSD является направленность зависимостей: компоненты нижнего слоя не могут зависеть от компонентов верхнего слоя. Это обеспечивает слабую связанность и высокую когезию модулей.

\subsection{Компонентная база shadcn/ui}

Для ускорения разработки интерфейса используется библиотека \textbf{shadcn/ui} \cite{shadcn2024} --- коллекция переиспользуемых React-компонентов, основанная на Radix UI primitives и стилизованная с помощью Tailwind CSS. Компоненты shadcn/ui копируются непосредственно в кодовую базу проекта, что обеспечивает полный контроль над их внешним видом и поведением.

\subsection{Управление состоянием и запросами}

Для управления серверным состоянием используется \textbf{TanStack Query} (React Query) \cite{tanstack2024} --- библиотека, предоставляющая декларативный подход к загрузке, кэшированию и обновлению данных с сервера. TanStack Query автоматизирует обработку состояний загрузки, ошибок и инвалидацию кэша.

Для глобального клиентского состояния используется библиотека \textbf{Zustand} --- лёгкое решение для управления состоянием с минимальным шаблонным кодом.


\section{Модуль генерации моделей}

\subsection{Конвейер генерации}

Модуль генерации реализует многоэтапный конвейер преобразования текстового описания системы в исполняемую системно-динамическую модель. Конвейер состоит из следующих последовательных этапов (рис.~\ref{fig:pipeline}):

\begin{figure}[H]
    \centering
    \begin{tikzpicture}[
        node distance=0.4cm,
        stage/.style={draw, rounded corners=4pt, minimum width=2.4cm, minimum height=1cm, align=center, font=\small},
        arrow/.style={-{Stealth[scale=0.8]}, thick}
    ]
        \node[stage, fill=blue!10] (text) {Текстовое\\описание};
        \node[stage, fill=green!10, right=0.6cm of text] (vars) {Переменные\\(LLM)};
        \node[stage, fill=green!10, right=0.6cm of vars] (cld) {CLD\\(LLM)};
        \node[stage, fill=green!10, right=0.6cm of cld] (eq) {Уравнения\\(LLM)};
        \node[stage, fill=green!10, right=0.6cm of eq] (sfd) {SFD\\(LLM/правила)};
        \node[stage, fill=yellow!15, right=0.6cm of sfd] (py) {Python\\(генератор)};
        \node[stage, fill=orange!10, right=0.6cm of py] (sim) {Симуляция\\(PySD/SciPy)};

        \draw[arrow] (text) -- (vars);
        \draw[arrow] (vars) -- (cld);
        \draw[arrow] (cld) -- (eq);
        \draw[arrow] (eq) -- (sfd);
        \draw[arrow] (sfd) -- (py);
        \draw[arrow] (py) -- (sim);
    \end{tikzpicture}
    \caption{Конвейер генерации системно-динамической модели (актуальная версия)}
    \label{fig:pipeline}
\end{figure}

\begin{enumerate}
    \item \textbf{Этап 1: Извлечение переменных.} На вход подаётся текстовое описание системы. LLM извлекает ключевые переменные с указанием типа (stock, flow, auxiliary, constant), единиц измерения, начальных значений и описания.

    \item \textbf{Этап 2: Генерация CLD.} На основе списка переменных LLM строит ориентированный граф причинно-следственных связей $G = (V, E)$, где $V$ --- множество переменных, $E \subseteq V \times V$ --- множество ориентированных рёбер с пометками полярности $(+, -)$. Граф валидируется на связность, наличие обратных связей и отсутствие изолированных вершин.

    \item \textbf{Этап 3: Генерация уравнений.} На основе переменных и CLD формируется система уравнений в структурированном JSON-формате. Каждое уравнение содержит левую часть (\texttt{lhs}), правую часть (\texttt{rhs}) и тип переменной. Уровни описываются через вызов \texttt{integ(чистый поток, начальное\_значение)}.

    \item \textbf{Этап 3b: Постпроцессинг уравнений.} Автоматически исправляются типичные ошибки LLM: константы и уровни с нулевыми или пустыми значениями заменяются на эвристически подобранные ненулевые defaults (например, 0.05 для коэффициентов, 100.0 для начальных запасов).

    \item \textbf{Этап 4: Проверка согласованности.} Проверяется соответствие между переменными, уравнениями и узлами CLD: каждая переменная имеет уравнение, каждое уравнение относится к известной переменной, все узлы CLD присутствуют в списке переменных.

    \item \textbf{Этап 5: Генерация SFD.} На основе переменных и уравнений строится диаграмма потоков и уровней. LLM формирует позиции узлов и связи, после чего выполняется валидация структуры SFD (наличие потоков для уровней, отсутствие изолированных элементов).

    \item \textbf{Этап 6: Детерминированная генерация Python-кода.} По структурированным уравнениям формируется исполняемый Python-код, совместимый с PySD. Все bare-ссылки на переменные преобразуются в вызовы функций от \texttt{time}, а функции \texttt{MAX}/\texttt{MIN}/\texttt{ABS} нормализуются к Python-встроенным \texttt{max}/\texttt{min}/\texttt{abs}. Данный этап не требует обращения к LLM, что делает генерацию воспроизводимой.

    \item \textbf{Этап 7: Генерация Vensim-кода.} Уравнения экспортируются в текстовый формат Vensim .mdl с преобразованием Python-функций (\texttt{np.exp}, \texttt{max} и др.) в Vensim-синтаксис.

    \item \textbf{Этап 8: Симуляция.} Python-код модели загружается в PySD; при ошибках, связанных с импортом \texttt{integ} или несовместимостью API, автоматически используется SciPy \texttt{solve\_ivp} fallback. Результаты симуляции (временные ряды) возвращаются пользователю.
\end{enumerate}

\subsection{Промпт-инжиниринг}

Качество генерации LLM критически зависит от структуры промптов. В проекте \texttt{sdCloud} используются \textbf{few-shot промпты}, включающие примеры входных данных и ожидаемых ответов для каждого этапа конвейера.

Структура промпта для генерации CLD:

\begin{lstlisting}[caption={Структура промпта для генерации CLD}]
SYSTEM: Ты --- эксперт по системной динамике. Твоя задача ---
построить диаграмму причинно-следственных связей (CLD) на основе
текстового описания системы.

Формат ответа: JSON со следующей структурой:
{
  "variables": ["var1", "var2", ...],
  "edges": [
    {"from": "var1", "to": "var2", "polarity": "+"},
    ...
  ],
  "loops": [
    {"type": "R", "variables": ["var1", "var2", "var3"]},
    ...
  ]
}

Пример:
Описание: "Популяция кроликов растет при наличии пищи.
Чем больше популяция, тем больше потребление пищи."

Ответ:
{
  "variables": ["Популяция кроликов", "Пищевые ресурсы",
               "Потребление пищи"],
  "edges": [
    {"from": "Пищевые ресурсы", "to": "Популяция кроликов",
     "polarity": "+"},
    {"from": "Популяция кроликов", "to": "Потребление пищи",
     "polarity": "+"},
    {"from": "Потребление пищи", "to": "Пищевые ресурсы",
     "polarity": "-"}
  ],
  "loops": [
    {"type": "R", "variables": ["Популяция кроликов",
     "Потребление пищи", "Пищевые ресурсы"]},
    {"type": "B", "variables": ["Популяция кроликов",
     "Потребление пищи", "Пищевые ресурсы"]}
  ]
}

Теперь построй CLD для следующего описания:
{user_description}
\end{lstlisting}

\subsection{Детерминированный генератор Python-кода}

Важным архитектурным решением, отличающим текущую версию платформы от ранних прототипов, является отказ от генерации исполняемого Python-кода непосредственно LLM. Вместо этого используется детерминированный генератор \texttt{python\_generator.py}, который строит код по структурированным уравнениям:

\begin{itemize}
    \item Каждая переменная превращается в функцию \texttt{def имя(time): return выражение}.
    \item Все bare-ссылки на переменные внутри правых частей автоматически преобразуются в вызовы функций, что устраняет проблему пропущенных скобок, характерную для LLM-генерации.
    \item Поддерживается нормализация Vensim-функций \texttt{MAX}, \texttt{MIN}, \texttt{ABS} в Python-встроенные аналоги.
    \item Результат воспроизводим: один и тот же набор уравнений всегда порождает идентичный код.
\end{itemize}

Этот подход существенно повышает надёжность симуляции, поскольку исключает целый класс ошибок, связанных с галлюцинациями LLM при генерации кода (неправильные импорты, неопределённые переменные, синтаксические ошибки).


\subsection{Fallback-симуляция через SciPy}

В процессе эксплуатации платформы выявлена несовместимость сгенерированного Python-кода (использующего \texttt{from pysd.py\_backend.functions import integ}) с актуальной версией PySD 3.14.x, в которой символ \texttt{integ} недоступен для импорта. Для сохранения работоспособности системы разработан fallback-модуль \texttt{scipy\_runner.py}:

\begin{itemize}
    \item Парсит сгенерированный Python-код через модуль \texttt{ast}.
    \item Выделяет уровни (stock) по наличию вызова \texttt{integ(поток, начальное\_значение)}.
    \item Формирует функцию \texttt{derivatives(t, state)}, вычисляющую чистые потоки для каждого уровня.
    \item Интегрирует систему ОДУ методами SciPy (RK45 по умолчанию, доступны LSODA, BDF, DOP853 и др.).
    \item Поддерживает переопределение параметров и начальных значений уровней.
\end{itemize}

Таким образом, симуляция остаётся работоспособной даже при отказе PySD, что повышает отказоустойчивость платформы.


\subsection{Форматы промежуточных представлений}

Промежуточные результаты каждого этапа конвейера сериализуются в формате JSON, обеспечивая прозрачность процесса и возможность отладки:

\textbf{Формат CLD} (этап 2):
\begin{lstlisting}[caption={Формат представления CLD}]
{
  "variables": ["Population", "Births", "Deaths",
                "Food", "Consumption"],
  "edges": [
    {"from": "Population", "to": "Births", "polarity": "+"},
    {"from": "Population", "to": "Deaths", "polarity": "+"},
    {"from": "Food", "to": "Births", "polarity": "+"},
    {"from": "Food", "to": "Consumption", "polarity": "+"},
    {"from": "Consumption", "to": "Food", "polarity": "-"}
  ],
  "loops": [
    {"type": "R", "variables": ["Population", "Births"]},
    {"type": "B", "variables": ["Population", "Deaths"]}
  ]
}
\end{lstlisting}

\textbf{Формат SFD} (этап 5):
\begin{lstlisting}[caption={Формат представления SFD}]
{
  "stocks": [
    {"name": "Population", "initial": 100,
     "inflows": ["Births"], "outflows": ["Deaths"]}
  ],
  "flows": [
    {"name": "Births", "formula": "Population * birth_rate"},
    {"name": "Deaths", "formula": "Population * death_rate"}
  ],
  "auxiliaries": [
    {"name": "birth_rate", "value": 0.05},
    {"name": "death_rate", "value": 0.02}
  ]
}
\end{lstlisting}

\textbf{Формат Vensim MDL} (этап 7):
\begin{lstlisting}[caption={Пример сгенерированного кода MDL}]
Population = INTEG(Births - Deaths, 100)
~ people
~

Births = Population * birth_rate
~ people/year
~

Deaths = Population * death_rate
~ people/year
~

birth_rate = 0.05
~ 1/year
~

death_rate = 0.02
~ 1/year
~
\end{lstlisting}


\section{Модуль валидации графов}

\subsection{Задачи валидации}

Модуль валидации графов выполняет проверку корректности топологии сгенерированных CLD и SFD перед переходом к генерации уравнений и симуляции. Валидация направлена на выявление типичных ошибок генерации LLM:

\begin{itemize}
    \item Переменные без входящих или исходящих связей (изолированные вершины).
    \item Некорректная классификация переменных (отсутствие уровней в SFD).
    \item Нарушение баланса потоков для уровней.
    \item Отсутствие замкнутых петель обратной связи.
    \item Ссылки на несуществующие переменные в уравнениях.
\end{itemize}

\subsection{Алгоритмы валидации}

Валидация реализована с использованием библиотеки \textbf{NetworkX} \cite{networkx2008} --- зрелой Python-библиотеки для создания, манипулирования и анализа графов.

\textbf{Проверка связности графа:}

Ориентированный граф CLD должен быть \emph{слабо связным} (weakly connected), то есть его неориентированная версия должна быть связным графом. Это условие гарантирует отсутствие изолированных подсистем:

\begin{algorithm}[H]
    \caption{Проверка слабой связности CLD}
    \begin{algorithmic}[1]
        \Require Ориентированный граф $G = (V, E)$
        \Ensure \text{True}, если граф слабо связен; \text{False} иначе
        \State $G_{undirected} \gets$ \Call{ToUndirected}{$G$}
        \State $components \gets$ \Call{ConnectedComponents}{$G_{undirected}$}
        \If{$|components| = 1$}
            \State \Return \text{True}
        \Else
            \State \Return \text{False}
        \EndIf
    \end{algorithmic}
\end{algorithm}

\textbf{Обнаружение петель обратной связи:}

Системно-динамическая модель должна содержать хотя бы одну петлю обратной связи. Для обнаружения петель используется алгоритм поиска элементарных циклов в ориентированном графе (алгоритм Джонсона \cite{johnson1975finding}):

\begin{algorithm}[H]
    \caption{Обнаружение петель обратной связи}
    \begin{algorithmic}[1]
        \Require Ориентированный граф $G = (V, E)$
        \Ensure Список циклов $C$
        \State $C \gets$ \Call{SimpleCycles}{$G$} \Comment{алгоритм Джонсона}
        \State $feedback\_loops \gets \emptyset$
        \For{$cycle$ \textbf{in} $C$}
            \State $sign \gets$ \Call{ComputeLoopSign}{$cycle, E$}
            \State \Call{Append}{$feedback\_loops, (cycle, sign)$}
        \EndFor
        \State \Return $feedback\_loops$
    \end{algorithmic}
\end{algorithm}

Знак петли определяется произведением полярностей её рёбер: чётное число отрицательных связей соответствует усиливающей петле ($R$), нечётное --- балансирующей ($B$).

\textbf{Проверка структуры SFD:}

Для SFD выполняется проверка корректности структурных элементов:

\begin{itemize}
    \item Каждый уровень имеет хотя бы один входной или выходной поток.
    \item Каждый поток связан ровно с одним уровнем (источником или стоком).
    \item Все переменные, упомянутые в формулах потоков, определены в модели.
\end{itemize}

\subsection{Результаты валидации}

Результаты проверки графа возвращаются в структурированном формате:

\begin{lstlisting}[caption={Формат результата валидации}]
{
  "valid": false,
  "errors": [
    {"type": "isolated_variable", "variable": "ExternalFactor"},
    {"type": "unbalanced_stock",
     "stock": "Inventory", "missing": "outflow"}
  ],
  "warnings": [
    {"type": "no_reinforcing_loop",
     "message": "Model has no reinforcing loops"}
  ],
  "stats": {
    "variables": 8,
    "edges": 12,
    "feedback_loops": 3,
    "reinforcing": 1,
    "balancing": 2
  }
}
\end{lstlisting}

При обнаружении ошибок модуль валидации формирует отчёт, который может быть использован для корректирующей генерации (iterative refinement) --- повторного запроса к LLM с указанием выявленных проблем.


\section{Интеграция с OpenRouter}

\subsection{Архитектура интеграции}

\textbf{OpenRouter} --- сервис-агрегатор, предоставляющий унифицированный REST API для доступа к различным LLM (GPT-4, Claude, LLaMA, Mistral, Gemma и др.). Взаимодействие с OpenRouter в платформе \texttt{sdCloud} реализовано через двухуровневый адаптер:

\begin{itemize}
    \item \textbf{UniversalAI} --- синхронный клиент, использующий библиотеку \texttt{litellm} и официальный OpenAI SDK для вызова чат-комплишенов, генерации изображений, аудио, видео и транскрипции.
    \item \textbf{OpenRouterLLM} --- асинхронная обёртка над \texttt{UniversalAI}, выполняющая блокирующие вызовы в отдельном потоке (\texttt{asyncio.to\_thread}) с заданным таймаутом, что позволяет использовать клиент в FastAPI без блокировки event loop.
\end{itemize}

Использование OpenRouter в качестве прослойки между платформой \texttt{sdCloud} и провайдерами LLM обеспечивает следующие преимущества:

\begin{itemize}
    \item \textbf{Унификация интерфейса} --- единый формат запросов и ответов для всех моделей через совместимый с OpenAI API.
    \item \textbf{Гибкость выбора модели} --- возможность переключения между различными LLM изменением только переменной окружения \texttt{OPENROUTER\_MODEL}.
    \item \textbf{Доступ к бесплатным моделям} --- возможность использования free-tier моделей (например, \texttt{openrouter/google/gemma-4-31b-it:free}) для прототипирования и тестирования.
    \item \textbf{Расширенные модальности} --- потенциал для генерации сопроводительных изображений, аудио- и видеоматериалов к моделям.
\end{itemize}

\subsection{Формат запросов}

Взаимодействие с OpenRouter API реализовано через HTTP-запросы к эндпоинту \texttt{/chat/completions}, совместимому с форматом OpenAI API. Пример запроса, формируемого через \texttt{litellm}:

\begin{lstlisting}[caption={Пример запроса к OpenRouter API}]
POST https://openrouter.ai/api/v1/chat/completions
Content-Type: application/json
Authorization: Bearer {OPENROUTER_API_KEY}

{
  "model": "openrouter/google/gemma-4-31b-it:free",
  "messages": [
    {"role": "system", "content": "system_prompt"},
    {"role": "user", "content": "user_description"}
  ],
  "temperature": 0.3,
  "max_tokens": 4096
}
\end{lstlisting}

Параметр \texttt{temperature} установлен на уровне 0.3--0.4 для снижения стохастичности генерации и повышения воспроизводимости результатов. Для извлечения переменных, CLD, уравнений и SFD ответ модели извлекается из markdown-блока \texttt{```json} и парсится как JSON; при неудаче выполняется fallback-поиск массива или объекта в тексте.

\subsection{Обработка ответов и ошибок}

Ответ от OpenRouter API содержит сгенерированный текст в поле \texttt{choices[0].message.content}. Модуль интеграции выполняет следующие операции:

\begin{enumerate}
    \item Извлечение JSON из markdown-блоков или прямой поиск JSON-объекта/массива в тексте.
    \item Валидация структуры ответа (наличие обязательных полей, корректность типов).
    \item Обработка ошибок API (таймауты, лимиты запросов, недоступность модели, ошибки аутентификации).
    \item Логирование запросов и ответов для отладки.
\end{enumerate}

При возникновении ошибок генерация конкретного артефакта не прерывает весь конвейер: пустые списки переменных, уравнений или рёбер заменяются на безопасные defaults, а информация об ошибке фиксируется в логе генерации.


\section{Гостевой режим работы}

В текущей версии платформы реализован гостевой режим, позволяющий анонимным пользователям знакомиться с возможностями \texttt{sdCloud} без создания учётной записи.

\subsection{Принципы гостевого режима}

\begin{itemize}
    \item \textbf{Statelessness} --- сервер не сохраняет данные гостевых пользователей в PostgreSQL.
    \item \textbf{История в браузере} --- полный контекст диалога и сгенерированные модели хранятся в памяти браузера (Zustand store).
    \item \textbf{Синхронная генерация} --- в отличие от авторизованного режима, гостевая генерация выполняется синхронно в рамках одного HTTP-запроса, что упрощает клиентскую логику.
    \item \textbf{Полный набор артефактов} --- гость получает переменные, CLD, уравнения, SFD, Python-код, Vensim-код и может запустить симуляцию.
\end{itemize}

\subsection{API гостевого режима}

Гостевые эндпоинты расположены по префиксу \texttt{/api/v1/guest}:

\begin{itemize}
    \item \texttt{POST /api/v1/guest/chat} --- отправка сообщения с синхронной генерацией модели и ответом ассистента.
    \item \texttt{POST /api/v1/guest/simulate} --- запуск симуляции по предоставленному Python- или Vensim-коду.
\end{itemize}

На frontend гостевой режим поддерживается отдельным слоем сущностей \texttt{entities/guest} и виджетом \texttt{chat-layout}, который корректно обрабатывает отсутствие авторизации и сохраняет историю в Zustand.


% ---------- Глава 3 ----------
\newpage
% ============================================
% Глава 3. Реализация и результаты
% ============================================

\chapter{РЕАЛИЗАЦИЯ И РЕЗУЛЬТАТЫ}

\section{Структура проекта}

\subsection{Организация кодовой базы}

Проект \texttt{sdCloud} организован в виде монорепозитория с разделением на backend и frontend части. Корневая структура проекта представлена на листинге~\ref{lst:project-tree}.

\begin{lstlisting}[caption={Структура директорий проекта sdCloud (актуальная версия)}, label={lst:project-tree}, language={}, basicstyle=\ttfamily\scriptsize]
sdcloud/
|-- backend/
|   |-- app/
|   |   |-- main.py              # Точка входа FastAPI
|   |   |-- api/
|   |   |   |-- v1/
|   |   |   |   |-- router.py    # Сборка маршрутов API v1
|   |   |   |   |-- endpoints/
|   |   |   |   |   |-- auth.py        # JWT-аутентификация
|   |   |   |   |   |-- conversations.py # CRUD диалогов
|   |   |   |   |   |-- generate.py    # Запуск генерации
|   |   |   |   |   |-- guest.py       # Гостевые эндпоинты
|   |   |   |   |   |-- health.py      # Health check
|   |   |   |   |   |-- models.py      # CRUD сгенерированных моделей
|   |   |   |   |   |-- simulate.py    # Эндпоинты симуляции
|   |   |-- core/
|   |   |   |-- config.py        # Pydantic Settings
|   |   |   |-- dependencies.py  # Зависимости FastAPI
|   |   |   |-- security.py      # JWT, bcrypt
|   |   |   |-- exceptions.py    # Кастомные исключения
|   |   |-- models/
|   |   |   |-- user.py          # Модель пользователя
|   |   |   |-- conversation.py  # Модель диалога
|   |   |   |-- message.py       # Модель сообщения
|   |   |   |-- generated_model.py # Модель SD-модели
|   |   |   |-- simulation.py    # Модель результата симуляции
|   |   |-- schemas/             # Pydantic-схемы запросов/ответов
|   |   |-- services/
|   |   |   |-- generation/
|   |   |   |   |-- pipeline.py        # Многоэтапный конвейер
|   |   |   |   |-- python_generator.py # Детерминированный генератор
|   |   |   |   |-- vensim_export.py   # Экспорт в .mdl
|   |   |   |-- graph/
|   |   |   |   |-- cld_builder.py     # Построение CLD
|   |   |   |   |-- sfd_builder.py     # Построение SFD
|   |   |   |   |-- validator.py       # Валидация графов
|   |   |   |-- llm/
|   |   |   |   |-- universal_ai.py    # LiteLLM + OpenAI SDK клиент
|   |   |   |   |-- openrouter.py      # Async-обертка
|   |   |   |-- simulation/
|   |   |   |   |-- pysd_runner.py     # Симуляция через PySD
|   |   |   |   |-- scipy_runner.py    # Fallback на SciPy solve_ivp
|   |   |-- utils/
|   |   |   |-- sd_prompts.py          # Шаблоны промптов
|   |   |   |-- model_postprocess.py   # Постпроцессинг уравнений
|   |   |-- workers/
|   |   |   |-- tasks.py               # Фоновые asyncio-задачи
|   |   |-- alembic/             # Миграции базы данных
|   |   |-- tests/               # Pytest-тесты
|   |-- pyproject.toml           # Зависимости (setuptools)
|   |-- Dockerfile
|   |-- requirements.txt         # Фиксированные зависимости
|
|-- frontend/
|   |-- src/
|   |   |-- app/                 # Next.js App Router: страницы
|   |   |-- widgets/             # Самостоятельные UI-блоки
|   |   |   |-- chat-layout/
|   |   |   |-- header/
|   |   |   |-- sidebar/
|   |   |-- features/            # Фичи: auth, chat-send, model-view,
|   |   |   |                    # simulation-run
|   |   |-- entities/            # Сущности: user, chat, model,
|   |   |   |                    # simulation, guest
|   |   |-- shared/              # API-клиент, UI-компоненты, конфиг
|   |-- package.json             # Зависимости (Bun)
|   |-- next.config.js
|   |-- tailwind.config.ts
|   |-- tsconfig.json
|
|-- docker-compose.yaml          # Общая оркестрация
|-- README.md
|-- .env.example
\end{lstlisting}

\subsection{Управление зависимостями}

Backend-проект использует стандартный \texttt{pip} с файлом \texttt{requirements.txt} (детерминированная установка в Docker) и \texttt{pyproject.toml} для метаданных проекта. В процессе развёртывания была выявлена и устранена проблема долгого разрешения зависимостей \texttt{litellm} путём фиксации версии \texttt{litellm==1.40.0}. Основные зависимости backend представлены в таблице~\ref{tab:backend-deps}.

\begin{table}[H]
    \centering
    \caption{Основные зависимости backend}
    \label{tab:backend-deps}
    \begin{tabular}{|l|l|l|}
        \hline
        \textbf{Пакет} & \textbf{Версия} & \textbf{Назначение} \\
        \hline
        fastapi & $\geq 0.110$ & Веб-фреймворк \\
        \hline
        uvicorn & $\geq 0.27$ & ASGI-сервер \\
        \hline
        sqlalchemy & $\geq 2.0$ & ORM \\
        \hline
        alembic & $\geq 1.13$ & Миграции БД \\
        \hline
        asyncpg & $\geq 0.29$ & Асинхронный драйвер PostgreSQL \\
        \hline
        redis & $\geq 5.0$ & Клиент Redis \\
        \hline
        taskiq & $\geq 0.11$ & Очередь задач \\
        \hline
        httpx & $\geq 0.27$ & Асинхронный HTTP-клиент \\
        \hline
        pydantic & $\geq 2.6$ & Валидация данных \\
        \hline
        pydantic-settings & $\geq 2.1$ & Управление конфигурацией \\
        \hline
        networkx & $\geq 3.2$ & Анализ графов \\
        \hline
        pysd & $\geq 3.13$ & Симуляция SD-моделей (попытка загрузки) \\
        \hline
        scipy & $\geq 1.11$ & Численное интегрирование fallback \\
        \hline
        litellm & $1.40.0$ & Унифицированный клиент LLM \\
        \hline
        openai & $\geq 1.0$ & OpenAI SDK для OpenRouter \\
        \hline
        pytest & $\geq 8.0$ & Фреймворк тестирования \\
        \hline
        pytest-asyncio & $\geq 0.23$ & Тестирование async-кода \\
        \hline
    \end{tabular}
\end{table}


\section{API эндпоинты}

\subsection{Структура REST API}

API платформы спроектировано в соответствии с принципами REST и использует формат JSON для обмена данными. Базовый путь API: \texttt{/api/v1/}. Основные эндпоинты системы представлены в таблице~\ref{tab:api-endpoints}.

\begin{table}[H]
    \centering
    \caption{Основные эндпоинты REST API}
    \label{tab:api-endpoints}
    \begin{tabular}{|p{4.5cm}|p{1.8cm}|p{7.5cm}|}
        \hline
        \textbf{Эндпоинт} & \textbf{Метод} & \textbf{Описание} \\
        \hline
        \texttt{/api/v1/health} & GET & Проверка состояния сервиса \\
        \hline
        \texttt{/api/v1/sessions} & GET & Получение списка сессий пользователя \\
        \hline
        \texttt{/api/v1/sessions} & POST & Создание новой сессии моделирования \\
        \hline
        \texttt{/api/v1/sessions/\{id\}} & GET & Получение сессии по идентификатору \\
        \hline
        \texttt{/api/v1/sessions/\{id\}} & DELETE & Удаление сессии \\
        \hline
        \texttt{/api/v1/sessions/\{id\}/messages} & GET & Получение сообщений сессии \\
        \hline
        \texttt{/api/v1/sessions/\{id\}/messages} & POST & Отправка сообщения (запуск генерации) \\
        \hline
        \texttt{/api/v1/models/\{id\}} & GET & Получение сгенерированной модели \\
        \hline
        \texttt{/api/v1/models/\{id\}/cld} & GET & Получение CLD модели \\
        \hline
        \texttt{/api/v1/models/\{id\}/sfd} & GET & Получение SFD модели \\
        \hline
        \texttt{/api/v1/models/\{id\}/mdl} & GET & Получение кода модели в формате MDL \\
        \hline
        \texttt{/api/v1/models/\{id\}/simulate} & POST & Запуск симуляции модели \\
        \hline
        \texttt{/api/v1/models/\{id\}/results} & GET & Получение результатов симуляции \\
        \hline
        \texttt{/api/v1/validate/cld} & POST & Валидация CLD \\
        \hline
        \texttt{/api/v1/validate/sfd} & POST & Валидация SFD \\
        \hline
        \texttt{/api/v1/guest/chat} & POST & Гостевое сообщение + генерация модели \\
        \hline
        \texttt{/api/v1/guest/simulate} & POST & Гостевая симуляция \\
        \hline
    \end{tabular}
\end{table}

\subsection{Примеры запросов и ответов}

\textbf{Создание сессии моделирования:}

\begin{lstlisting}[caption={Запрос на создание сессии}, language={}]
POST /api/v1/sessions
Content-Type: application/json

{
  "title": "Модель предприятия",
  "description": "Модель роста и падения производства"
}
\end{lstlisting}

\begin{lstlisting}[caption={Ответ при создании сессии}, language={}]
{
  "id": "a1b2c3d4-e5f6-7890-abcd-ef1234567890",
  "title": "Модель предприятия",
  "description": "Модель роста и падения производства",
  "created_at": "2025-03-15T10:30:00Z",
  "updated_at": "2025-03-15T10:30:00Z",
  "status": "active"
}
\end{lstlisting}

\textbf{Отправка сообщения и запуск генерации:}

\begin{lstlisting}[caption={Запрос на генерацию модели}, language={}]
POST /api/v1/sessions/a1b2c3d4-e5f6-7890-abcd-ef1234567890/messages
Content-Type: application/json

{
  "content": "Предприятие нанимает сотрудников для увеличения
              производительности. Рост производительности
              приводит к увеличению выручки. Выручка определяет
              бюджет на найм. Рост числа сотрудников замедляется
              из-за ограниченности офисных помещений.",
  "generate_model": true
}
\end{lstlisting}

\begin{lstlisting}[caption={Ответ при запуске генерации}, language={}]
{
  "message_id": "msg-12345",
  "task_id": "task-67890",
  "status": "queued",
  "estimated_time": "30-60 seconds"
}
\end{lstlisting}


\section{Модели данных}

\subsection{Схема базы данных}

Реляционная схема базы данных PostgreSQL включает следующие основные таблицы (рис.~\ref{fig:er-diagram}):

\begin{figure}[H]
    \centering
    \begin{tikzpicture}[
        entity/.style={draw, rectangle, minimum width=3.5cm, minimum height=0.8cm, align=center, font=\small},
        relation/.style={draw, diamond, aspect=2, align=center, font=\small},
        arrow/.style={-Stealth, thick}
    ]
        \node[entity, fill=blue!10] (user) {\textbf{users}\\id, email, name, ...};
        \node[entity, fill=blue!10, right=3cm of user] (session) {\textbf{sessions}\\id, title, status, ...};
        \node[entity, fill=blue!10, right=3cm of session] (message) {\textbf{messages}\\id, content, role, ...};
        \node[entity, fill=green!10, below=2cm of session] (model) {\textbf{sd\_models}\\id, clod\_json, sfd\_json, ...};
        \node[entity, fill=green!10, below=2cm of message] (sim) {\textbf{simulations}\\id, results, params, ...};

        \draw[arrow] (user) -- (session) node[midway, above, font=\scriptsize] {1:N};
        \draw[arrow] (session) -- (message) node[midway, above, font=\scriptsize] {1:N};
        \draw[arrow] (message) -- (model) node[midway, left, font=\scriptsize] {1:1};
        \draw[arrow] (model) -- (sim) node[midway, above, font=\scriptsize] {1:N};
    \end{tikzpicture}
    \caption{ER-диаграмма базы данных sdCloud}
    \label{fig:er-diagram}
\end{figure}

\subsection{Описание таблиц}

\textbf{Таблица \texttt{users}} хранит информацию о пользователях платформы:

\begin{table}[H]
    \centering
    \caption{Структура таблицы \texttt{users}}
    \begin{tabular}{|l|l|l|}
        \hline
        \textbf{Поле} & \textbf{Тип} & \textbf{Описание} \\
        \hline
        id & UUID (PK) & Уникальный идентификатор \\
        \hline
        email & VARCHAR(255) & Email пользователя (уникальный) \\
        \hline
        hashed\_password & VARCHAR(255) & Хеш пароля (bcrypt) \\
        \hline
        full\_name & VARCHAR(255) & Полное имя \\
        \hline
        is\_active & BOOLEAN & Флаг активности \\
        \hline
        created\_at & TIMESTAMP & Время создания \\
        \hline
    \end{tabular}
\end{table}

\textbf{Таблица \texttt{sessions}} представляет сессии моделирования:

\begin{table}[H]
    \centering
    \caption{Структура таблицы \texttt{sessions}}
    \begin{tabular}{|l|l|l|}
        \hline
        \textbf{Поле} & \textbf{Тип} & \textbf{Описание} \\
        \hline
        id & UUID (PK) & Уникальный идентификатор \\
        \hline
        user\_id & UUID (FK) & Ссылка на пользователя \\
        \hline
        title & VARCHAR(255) & Название сессии \\
        \hline
        description & TEXT & Описание \\
        \hline
        status & VARCHAR(50) & Статус (active, archived) \\
        \hline
        created\_at & TIMESTAMP & Время создания \\
        \hline
        updated\_at & TIMESTAMP & Время обновления \\
        \hline
    \end{tabular}
\end{table}

\textbf{Таблица \texttt{messages}} хранит сообщения в рамках сессий:

\begin{table}[H]
    \centering
    \caption{Структура таблицы \texttt{messages}}
    \begin{tabular}{|l|l|l|}
        \hline
        \textbf{Поле} & \textbf{Тип} & \textbf{Описание} \\
        \hline
        id & UUID (PK) & Уникальный идентификатор \\
        \hline
        session\_id & UUID (FK) & Ссылка на сессию \\
        \hline
        role & VARCHAR(20) & Роль (user, assistant, system) \\
        \hline
        content & TEXT & Содержание сообщения \\
        \hline
        model\_id & UUID (FK, nullable) & Ссылка на сгенерированную модель \\
        \hline
        task\_id & VARCHAR(100) & Идентификатор задачи Taskiq \\
        \hline
        created\_at & TIMESTAMP & Время создания \\
        \hline
    \end{tabular}
\end{table}

\textbf{Таблица \texttt{sd\_models}} хранит сгенерированные системно-динамические модели:

\begin{table}[H]
    \centering
    \caption{Структура таблицы \texttt{sd\_models}}
    \begin{tabular}{|l|l|l|}
        \hline
        \textbf{Поле} & \textbf{Тип} & \textbf{Описание} \\
        \hline
        id & UUID (PK) & Уникальный идентификатор \\
        \hline
        message\_id & UUID (FK) & Ссылка на сообщение \\
        \hline
        clod\_json & JSONB & CLD в формате JSON \\
        \hline
        sfd\_json & JSONB & SFD в формате JSON \\
        \hline
        mdl\_code & TEXT & Код модели в формате MDL \\
        \hline
        validation\_result & JSONB & Результат валидации \\
        \hline
        status & VARCHAR(50) & Статус (pending, valid, invalid) \\
        \hline
        created\_at & TIMESTAMP & Время создания \\
        \hline
    \end{tabular}
\end{table}

Использование типа \texttt{JSONB} для хранения структурированных данных (CLD, SFD, результаты валидации) обеспечивает гибкость схемы при сохранении возможности индексации и эффективного запроса вложенных полей.


\section{Детерминированный генератор Python-кода}

Ключевым элементом надёжности текущей версии является детерминированный генератор \texttt{python\_generator.py}. Он преобразует структурированные уравнения, полученные от LLM, в исполняемый Python-код, не прибегая повторно к языковой модели.

\subsection{Принцип работы}

Генератор выполняет следующие операции:

\begin{enumerate}
    \item Собирает множество имён переменных модели.
    \item Для каждого уравнения парсит правую часть с помощью \texttt{ast.parse}.
    \item Преобразует bare-ссылки на переменные в вызовы функций от \texttt{time}: \texttt{Population} становится \texttt{Population(time)}.
    \item Нормализует Vensim-функции \texttt{MAX}, \texttt{MIN}, \texttt{ABS} к Python-встроенным.
    \item Формирует модуль с импортами \texttt{numpy} и \texttt{integ} из PySD.
\end{enumerate}

\begin{lstlisting}[caption={Детерминированная генерация Python-кода (упрощённо)}, language=Python, basicstyle=\ttfamily\scriptsize]
def generate_python_code(variables, equations):
    variable_names = {v["name"] for v in variables}
    transformer = _build_call_transformer(variable_names)
    lines = ["import numpy as np",
             "from pysd.py_backend.functions import integ"]

    for eq in equations:
        lhs = eq.get("lhs", "")
        rhs = str(eq.get("rhs", "")).strip()
        rhs_tree = ast.parse(rhs, mode="eval")
        rhs_tree = transformer.visit(rhs_tree)
        transformed_rhs = ast.unparse(rhs_tree.body)
        lines.append(f"\ndef {lhs}(time):")
        lines.append(f"    return {transformed_rhs}")

    return "\n".join(lines)
\end{lstlisting}

Преимущество этого подхода --- воспроизводимость и отсутствие типичных для LLM ошибок: неправильных импортов, ссылок на несуществующие переменные, пропущенных скобок.


\section{Постпроцессинг уравнений}

LLM часто генерируют нулевые или пустые значения для констант и начальные значения уровней, равные нулю, что приводит к вырожденным симуляциям. Модуль \texttt{model\_postprocess.py} автоматически исправляет такие случаи эвристически подобранными defaults.

\begin{lstlisting}[caption={Постпроцессинг уравнений (упрощённо)}, language=Python, basicstyle=\ttfamily\scriptsize]
def post_process_equations(variables, equations):
    patched = []
    for eq in equations:
        eq = dict(eq)
        vtype = eq.get("variable_type", "").lower()
        if vtype == "constant":
            eq = _patch_constant_equation(eq, variables)
        elif vtype == "stock":
            eq = _patch_stock_equation(eq, variables)
        patched.append(eq)
    return patched
\end{lstlisting}

Для констант применяются name-based эвристики: коэффициенты получают 0.05, ставки роста --- 0.02, ёмкости --- 1000.0. Для уровней начальное значение берётся из переменной, если оно ненулевое, иначе устанавливается 100.0.


\section{Fallback-симуляция через SciPy}

При попытке загрузки сгенерированного кода в PySD 3.14.x возникает ошибка импорта \texttt{integ}. Для сохранения работоспособности разработан fallback на \texttt{scipy.integrate.solve\_ivp}.

\subsection{Алгоритм работы}

\begin{enumerate}
    \item Парсинг Python-кода через \texttt{ast}; выделение функций-уровней по вызову \texttt{integ(поток, начальное\_значение)}.
    \item Построение функции \texttt{derivatives(t, state)}, которая обновляет shared-состояние, вычисляет алгебраические переменные и возвращает чистые потоки для каждого уровня.
    \item Формирование начального состояния \texttt{initial\_state()}.
    \item Интегрирование системы ОДУ методом RK45 (или LSODA, BDF, DOP853) на равномерной сетке времени.
\end{enumerate}

\begin{lstlisting}[caption={Запуск симуляции через SciPy fallback}, language=Python, basicstyle=\ttfamily\scriptsize]
def run_simulation(python_code, params=None,
                   time_start=0.0, time_end=100.0,
                   time_step=0.25, integration_method="rk45"):
    namespace, stock_names, initial_values = \
        _build_simulation_module(python_code)

    # Переопределение параметров
    for param_name, param_value in (params or {}).items():
        if param_name in stock_names:
            # начальные значения уровней
            ...
        elif param_name in namespace:
            namespace[param_name] = lambda t, v=param_value: v

    derivatives_fn = namespace["derivatives"]
    initial_state = np.array(namespace["initial_state"](), dtype=float)

    n_steps = int((time_end - time_start) / time_step) + 1
    t_eval = np.linspace(time_start, time_end, n_steps)

    sol = solve_ivp(derivatives_fn,
                    [time_start, time_end],
                    initial_state,
                    method=method_map[integration_method],
                    t_eval=t_eval)

    return {"time": sol.t.tolist(),
            **{name: sol.y[i].tolist()
               for i, name in enumerate(stock_names)}}
\end{lstlisting}


\section{Процесс генерации моделей}

\subsection{Асинхронный конвейер}

Для авторизованных пользователей процесс генерации модели запускается как фоновая \texttt{asyncio}-задача через \texttt{asyncio.create\_task}, что обеспечивает:

\begin{itemize}
    \item Неблокирующую обработку запросов (пользователь сразу получает ответ ассистента).
    \item Устойчивость к таймаутам HTTP (генерация может занимать 30--120 секунд).
    \item Возможность дальнейшего перехода на полноценную очередь задач (Taskiq, Celery и др.).
\end{itemize}

В гостевом режиме генерация выполняется синхронно в рамках одного запроса, поскольку не требует записи в базу данных.

\subsection{Пошаговая реализация конвейера}

Конвейер реализован в функции \texttt{generate\_model\_from\_description}:

\begin{lstlisting}[caption={Основная функция конвейера генерации}, language=Python, basicstyle=\ttfamily\scriptsize]
async def generate_model_from_description(description: str) -> dict:
    llm = OpenRouterLLM()
    validator = GraphValidator()
    cld_builder = CLDBuilder(llm=llm)
    sfd_builder = SFDBuilder(llm=llm)
    exporter = VensimExporter()

    # Stage 1: переменные
    variables = await llm.chat_with_json_output(
        prompt=VARIABLE_EXTRACTION_PROMPT.format(
            description=description))

    # Stage 2: CLD
    cld_result = await cld_builder.build_from_variables(variables)

    # Stage 3: уравнения
    equations = await llm.chat_with_json_output(
        prompt=EQUATION_GENERATION_PROMPT.format(
            variables=json.dumps(variables),
            cld=json.dumps(cld_result)))

    # Stage 3b: постпроцессинг defaults
    equations = post_process_equations(variables, equations)

    # Stage 4: согласованность
    consistency = validator.validate_model_consistency(
        variables, equations,
        cld_result["nodes"], cld_result["edges"])

    # Stage 5: SFD
    sfd_result = await sfd_builder.build_from_equations(
        variables, equations)

    # Stage 6: Python-код (детерминированно)
    python_code = generate_python_code(variables, equations)

    # Stage 7: Vensim-экспорт
    vensim_code = exporter.export_to_vensim(variables, equations)

    return {
        "status": "completed",
        "variables": variables,
        "equations": equations,
        "cld_data": { ... },
        "sfd_data": { ... },
        "python_code": python_code,
        "vensim_code": vensim_code,
        "log": result_log,
    }
\end{lstlisting}


\section{Тестирование}

\subsection{Стратегия тестирования}

Тестирование платформы \texttt{sdCloud} проводится на нескольких уровнях:

\begin{enumerate}
    \item \textbf{Модульные тесты} (unit tests) --- проверка отдельных функций и классов.
    \item \textbf{Интеграционные тесты} --- проверка взаимодействия компонентов.
    \item \textbf{API-тесты} --- проверка эндпоинтов REST API.
    \item \textbf{End-to-end тесты} --- проверка полного сценария использования.
\end{enumerate}

\subsection{Модульные тесты}

Модульные тесты реализованы с использованием фреймворка \textbf{pytest} и его расширения \textbf{pytest-asyncio} для тестирования асинхронного кода. Покрытие тестами включает следующие модули:

\begin{lstlisting}[caption={Пример модульного теста валидации CLD}, language=Python, basicstyle=\ttfamily\scriptsize]
import pytest
import networkx as nx
from app.services.validate import validate_cld
from app.schemas.clod import CLDSchema, Edge

@pytest.mark.parametrize("cld,expected_valid", [
    # Корректный CLD с обратной связью
    (CLDSchema(
        variables=["A", "B", "C"],
        edges=[
            Edge(from_="A", to="B", polarity="+"),
            Edge(from_="B", to="C", polarity="+"),
            Edge(from_="C", to="A", polarity="-")
        ],
        loops=[{"type": "B", "variables": ["A", "B", "C"]}]
    ), True),
    # CLD без обратных связей
    (CLDSchema(
        variables=["A", "B"],
        edges=[Edge(from_="A", to="B", polarity="+")],
        loops=[]
    ), False),
    # CLD с изолированной переменной
    (CLDSchema(
        variables=["A", "B", "C"],
        edges=[
            Edge(from_="A", to="B", polarity="+")
        ],
        loops=[]
    ), False),
])
def test_validate_cld(cld, expected_valid):
    result = validate_cld(cld)
    assert result.valid == expected_valid
\end{lstlisting}

\subsection{Интеграционные тесты}

Интеграционные тесты проверяют взаимодействие компонентов с использованием тестовой базы данных (SQLite в памяти) и моков внешних сервисов:

\begin{lstlisting}[caption={Пример интеграционного теста API}, language=Python, basicstyle=\ttfamily\scriptsize]
import pytest
from httpx import AsyncClient
from app.main import app

@pytest.mark.asyncio
async def test_create_session(async_client: AsyncClient):
    response = await async_client.post("/api/v1/sessions", json={
        "title": "Test Session",
        "description": "Test description"
    })
    assert response.status_code == 201
    data = response.json()
    assert data["title"] == "Test Session"
    assert "id" in data

@pytest.mark.asyncio
async def test_generate_model_flow(async_client: AsyncClient,
                                   mock_llm_client):
    # Создание сессии
    session = await async_client.post("/api/v1/sessions", json={
        "title": "Population Model"
    })
    session_id = session.json()["id"]

    # Отправка запроса на генерацию
    response = await async_client.post(
        f"/api/v1/sessions/{session_id}/messages",
        json={"content": "Population growth model",
              "generate_model": True}
    )
    assert response.status_code == 202
    assert "task_id" in response.json()
\end{lstlisting}

\subsection{Тестовые данные}

В ходе разработки платформы была выполнена серия практических прогонов на типовых задачах системной динамики, включая модель роста популяции кроликов с ограниченными ресурсами и модель динамики рынка с функциями \texttt{MAX}. В процессе были выявлены и устранены следующие проблемы:

\begin{itemize}
    \item \textbf{Ошибка импорта \texttt{integ} в PySD 3.14.x} --- решена fallback-симуляцией через SciPy \texttt{solve\_ivp}.
    \item \textbf{Нулевые defaults у констант и уровней} --- решена эвристической подстановкой realistic defaults в постпроцессоре.
    \item \textbf{Недетерминированная генерация Python-кода LLM} --- заменена детерминированным генератором по уравнениям.
    \item \textbf{Функции \texttt{MAX}/\texttt{MIN}/\texttt{ABS} в верхнем регистре} --- нормализованы к Python-встроенным аналогам.
    \item \textbf{Тёмная тема визуализации графов} --- улучшена контрастность узлов и рёбер React Flow.
    \item \textbf{Проблемы авторизации в гостевом режиме} --- устранены редиректы 401 для анонимных пользователей.
\end{itemize}

Для тестирования генерации используется набор тестовых описаний систем различной сложности:

\begin{table}[H]
    \centering
    \caption{Тестовые кейсы для генерации моделей}
    \label{tab:test-cases}
    \begin{tabular}{|p{1.8cm}|p{10cm}|p{2cm}|}
        \hline
        \textbf{Код} & \textbf{Описание} & \textbf{Сложность} \\
        \hline
        TC-01 & Модель роста популяции с ограниченными ресурсами & Низкая \\
        \hline
        TC-02 & Модель эпидемии SIR (восприимчивые, инфицированные, выздоровевшие) & Низкая \\
        \hline
        TC-03 & Модель производственного предприятия с цепью поставок & Средняя \\
        \hline
        TC-04 & Модель динамики цен на рынке с конкуренцией & Средняя \\
        \hline
        TC-05 & Модель городской мобильности с несколькими видами транспорта & Высокая \\
        \hline
        TC-06 & Модель распространения инноваций (adoption model) & Низкая \\
        \hline
        TC-07 & Модель управления запасами с задержками & Средняя \\
        \hline
    \end{tabular}
\end{table}

\subsection{Результаты тестирования}

Результаты тестирования показали, что платформа успешно обрабатывает запросы генерации и симуляции для тестовых кейсов низкой и средней сложности. Детерминированный генератор Python-кода и fallback-симуляция через SciPy существенно повысили стабильность конвейера: типовые модели (рост популяции, рыночная динамика) успешно проходят генерацию и запускаются на симуляцию.

Для кейсов высокой сложности и при использовании бесплатных 30B-моделей на OpenRouter наблюдается снижение качества генерации: нестабильный структурированный вывод, ограниченные рассуждения, возможные отказы API. В текущей версии для повышения надёжности применяются жёсткий постпроцессинг, fallback-механизмы и детерминированная генерация кода. Дальнейшее развитие в сторону мультиагентной системы и сложных сценариев требует более мощных и стабильных языковых моделей.


% ---------- Глава 4 ----------
\newpage
% ============================================
% Глава 4. Обсуждение и направления развития
% ============================================

\chapter{ОБСУЖДЕНИЕ И НАПРАВЛЕНИЯ РАЗВИТИЯ}

\section{Ограничения текущей реализации}

\subsection{Ограничения генеративных возможностей}

Текущая реализация платформы \texttt{sdCloud} обладает рядом ограничений, определяющих границы применимости системы и направления дальнейшего совершенствования.

\textbf{Ограничения больших языковых моделей.} Качество генерируемых моделей напрямую зависит от возможностей используемой LLM. В текущей версии платформы используются в основном бесплатные модели семейства Gemma (30B параметров) на OpenRouter, что накладывает дополнительные ограничения:

\begin{itemize}
    \item \emph{Галлюцинации и фактические ошибки.} LLM могут генерировать структурно корректные, но содержательно ошибочные причинно-следственные связи, нулевые значения констант, пропускать скобки в формулах или использовать неподдерживаемые функции (например, \texttt{MAX} в верхнем регистре).

    \item \emph{Ограничение контекста и внимания.} Даже модели с формально большим контекстом испытывают трудности при генерации моделей с числом переменных более 10--15. Сложные системы требуют декомпозиции или более мощных моделей.

    \item \emph{Недетерминированность и нестабильность free-tier.} Идентичный запрос может давать различные результаты; бесплатные модели могут быть временно недоступны, иметь низкие rate limits и отказывать в генерации. Это затрудняет воспроизводимость и надёжность системы.

    \item \emph{Ограниченные рассуждения и structured output.} 30B-модели хуже справляются с многошаговыми рассуждениями, критикой собственных ответов и строгим соблюдением JSON-схем, что критично для мультиагентных сценариев.

    \item \emph{Ограниченные знания предметных областей.} LLM обладают широкими, но поверхностными знаниями. Для специализированных предметных областей качество генерации снижается.
\end{itemize}

В текущей версии влияние перечисленных ограничений частично компенсируется детерминированным генератором Python-кода, постпроцессингом уравнений, fallback-симуляцией и валидацией графов.

\textbf{Ограничения конвейера генерации и визуализации.} Многоэтапный конвейер, обеспечивающий прозрачность процесса, одновременно вносит дополнительные точки отказа:

\begin{itemize}
    \item Ошибки на ранних этапах (переменные, CLD) распространяются на последующие этапы (SFD, уравнения), поскольку каждый этап использует результат предыдущего.
    \item Межэтапная валидация обнаруживает ошибки, но механизм автоматической коррекции (iterative refinement) требует дополнительных итераций взаимодействия с LLM, что увеличивает время генерации и стоимость.
    \item Конвейер не поддерживает параллельную генерацию альтернативных вариантов модели для одного описания.
    \item Визуализация SFD и CLD на базе React Flow использует упрощённую ручную/круговую раскладку; при большом числе узлов рёбра могут пересекаться, а текст узлов остаётся мелким.
    \item Отсутствует автоматическая калибровка параметров по историческим данным и полноценный анализ чувствительности.
\end{itemize}

\subsection{Ограничения программной реализации}

\textbf{Масштабируемость.} Текущая архитектура предполагает развёртывание на одном сервере. При увеличении нагрузки необходимо внедрение горизонтального масштабирования (кластеризация воркеров Taskiq, репликация PostgreSQL, шардирование Redis).

\textbf{Аутентификация и авторизация.} В текущей реализации используется упрощённая система аутентификации на базе JWT-токенов. Для промышленного использования требуется интеграция с внешними провайдерами идентификации (OAuth 2.0, SAML), системой ролей и разграничением доступа к сессиям и моделям.

\textbf{Обработка ошибок.} Механизм обработки ошибок генерации требует доработки: необходимо реализовать более информативные сообщения об ошибках, систему логирования с разделением по уровням и мониторинг состояния системы.

\textbf{Совместимость форматов.} Текущая реализация поддерживает экспорт моделей в Python-код (основной исполняемый формат) и в текстовый формат Vensim .mdl. Для полноценной интеграции с экосистемой инструментов требуется поддержка формата XMILE, корректная работа с PySD 3.x и возможность импорта/экспота Stella.


\section{Планируемые улучшения}

\subsection{Краткосрочные улучшения текущей версии}

В рамках текущей версии платформы (до перехода к мультиагентной архитектуре) планируется реализация следующих улучшений:

\begin{itemize}
    \item \textbf{Улучшение визуализации графов.} Интеграция библиотек автоматической раскладки (\texttt{dagre}, \texttt{elk.js}, \texttt{cola}) для уменьшения пересечений рёбер в CLD/SFD и улучшения читаемости узлов.
    \item \textbf{Редактор моделей.} Возможность ручного редактирования переменных, уравнений и графов в интерфейсе с немедленным пересчётом Python-кода и симуляции.
    \item \textbf{Полноценная интеграция с PySD 3.x.} Адаптация генератора Python-кода под актуальный API PySD либо полный переход на SciPy-раннер с поддержкой табличных функций, задержек и lookup-функций.
    \item \textbf{Анализ чувствительности.} Автоматический запуск симуляций с вариацией ключевых параметров и визуализация диапазонов поведения модели.
    \item \textbf{Валидация уравнений.} Символьная проверка размерностей, поиск циклов в уравнениях, проверка корректности функций \texttt{integ}.
\end{itemize}

\subsection{Интеграция локальных и платных языковых моделей}

Приоритетным направлением развития является расширение поддержки моделей:

\begin{itemize}
    \item \textbf{Платные модели} (Claude 3.5 Sonnet, GPT-4o-mini, Gemini Flash) для повышения качества сложной генерации, критики и рассуждений.
    \item \textbf{Локальные LLM} через \textbf{Ollama} или \textbf{vLLM} для обеспечения конфиденциальности, offline-режима и снижения операционных затрат.
\end{itemize}

Использование free-tier 30B-моделей в текущей версии является компромиссом для прототипирования; для серьёзной мультиагентной системы и сложных моделей требуются более надёжные модели.

\subsection{Отложенное направление: мультиагентная система}

Разработка полноценной мультиагентной системы (агенты извлечения, построения CLD, генерации уравнений, критика, симуляции и др.) отложена до момента обеспечения:

\begin{itemize}
    \item Надёжного structured output и function calling используемой LLM.
    \item Достаточного контекстного окна для передачи состояния между агентами.
    \item Стабильной работы API без частых rate limits и отказов.
\end{itemize}

В текущей версии функции потенциальных агентов частично реализованы в виде отдельных модулей (\texttt{CLDBuilder}, \texttt{SFDBuilder}, \texttt{GraphValidator}, \texttt{python\_generator.py}), что создаёт основу для будущего перехода к мультиагентной архитектуре.

\subsection{Коллаборативное моделирование и интеграция с данными}

Перспективными направлениями являются:

\begin{itemize}
    \item Совместное редактирование моделей несколькими пользователями в реальном времени.
    \item Версионирование моделей и система комментариев.
    \item Загрузка исторических данных для калибровки параметров и валидации прогнозов.
    \item Публикация моделей в открытый репозиторий.
\end{itemize}


\section{Потенциальные приложения}

\subsection{Образование}

Платформа \texttt{sdCloud} имеет значительный потенциал в образовательной сфере:

\begin{itemize}
    \item \textbf{Обучение системной динамике.} Студенты могут быстро создавать модели на основе текстовых описаний, фокусируясь на анализе структуры и поведения системы, а не на технических аспектах моделирования.
    \item \textbf{Интерактивные лекции.} Преподаватели могут демонстрировать динамику различных систем в реальном времени, просто описывая их словами.
    \item \textbf{Проверка понимания.} Система может генерировать модели по описаниям студентов, позволяя преподавателю оценить глубину понимания предмета.
    \item \textbf{Междисциплинарное обучение.} Специалисты из различных областей (экология, экономика, медицина) могут создавать модели без глубокой подготовки в системной динамике.
\end{itemize}

\subsection{Бизнес-аналитика}

В сфере бизнеса платформа может применяться для:

\begin{itemize}
    \item \textbf{Быстрого прототипирования моделей.} Бизнес-аналитики могут оперативно создавать модели бизнес-процессов для оценки последствий решений.
    \item \textbf{Стратегического планирования.} Моделирование сценариев развития бизнеса с учётом обратных связей и запаздываний.
    \item \textbf{Управление цепями поставок.} Анализ динамики запасов, производства и спроса с выявлением эффектов хлыстовой верёвки (bullwhip effect).
\end{itemize}

\subsection{Анализ государственной политики}

Системная динамика традиционно применяется для анализа последствий государственных решений. Платформа \texttt{sdCloud} может снизить барьер для использования методологии в государственных учреждениях:

\begin{itemize}
    \item Моделирование динамики населения и демографических процессов.
    \item Анализ эпидемиологических сценариев и оценка эффективности вмешательств.
    \item Оценка последствий экологической политики и перехода к устойчивому развитию.
    \item Моделирование динамики транспортных систем и планирование инфраструктуры.
\end{itemize}

\subsection{Научные исследования}

В академической среде платформа может служить инструментом:

\begin{itemize}
    \item Быстрой проверки гипотез о структуре сложных систем.
    \item Автоматизации рутинных этапов моделирования для фокусировки на содержательном анализе.
    \item Создания репозиториев моделей с возможностью воспроизведения и верификации.
    \item Интеграции моделей системной динамики с методами машинного обучения (гибридное моделирование).
\end{itemize}

\subsection{Разработка программного обеспечения}

Методологии, отработанные в проекте \texttt{sdCloud}, имеют потенциал для применения в смежных областях:

\begin{itemize}
    \item Генерация формальных спецификаций из естественно-языковых требований.
    \item Автоматическое построение UML-диаграмм из текстовых описаний систем.
    \item Создание цифровых двойников на основе текстовой документации.
\end{itemize}

Таким образом, результаты настоящего исследования имеют междисциплинарную значимость и могут служить основой для дальнейших разработок на стыке искусственного интеллекта, системной динамики и программной инженерии.


% ---------- Заключение ----------
\newpage
% ============================================
% Заключение
% ============================================

\chapter*{ЗАКЛЮЧЕНИЕ}
\addcontentsline{toc}{chapter}{ЗАКЛЮЧЕНИЕ}

\section*{Итоги работы}
\addcontentsline{toc}{section}{Итоги работы}

В ходе выполнения настоящей научно-исследовательской работы был проведён комплексный анализ предметной области на стыке системной динамики и больших языковых моделей, спроектирована и реализована программная платформа \texttt{sdCloud}, предназначенная для автоматизированной генерации системно-динамических моделей на основе текстовых описаний.

В рамках работы были достигнуты следующие результаты:

\begin{enumerate}
    \item \textbf{Обзор предметной области.} Проведён систематический анализ теоретических основ системной динамики, включая концептуальное моделирование (CLD), построение диаграмм потоков и уровней (SFD) и формализацию уравнений. Исследованы существующие программные платформы моделирования (Vensim, Stella, AnyLogic, PySD), технологии больших языковых моделей (GPT-4, Claude, семейства LLaMA и Mistral) и существующие работы по интеграции LLM с системной динамикой (ChatPySD, исследования Akhavan, Liu и др.).

    \item \textbf{Проектирование архитектуры.} Разработана архитектура платформы \texttt{sdCloud}, включающая backend на базе FastAPI с использованием SQLAlchemy ORM, PostgreSQL, Redis и асинхронной очереди задач Taskiq; frontend на базе Next.js с применением методологии Feature-Sliced Design; а также модули генерации моделей, валидации графов и интеграции с OpenRouter API.

    \item \textbf{Разработка конвейера генерации.} Предложен и реализован многоэтапный конвейер генерации системно-динамических моделей, включающий извлечение переменных, построение CLD, генерацию уравнений, построение SFD, детерминированную генерацию Python-кода и Vensim-экспорт. Конвейер использует few-shot промптинг для структурированной генерации данных LLM.

    \item \textbf{Детерминированный генератор Python-кода.} Разработан модуль \texttt{python\_generator.py}, формирующий исполняемый код по структурированным уравнениям без повторного обращения к LLM, что повышает воспроизводимость и снижает влияние галлюцинаций.

    \item \textbf{Fallback-симуляция через SciPy.} Реализован модуль \texttt{scipy\_runner.py}, обеспечивающий симуляцию сгенерированных моделей методами \texttt{solve\_ivp} при невозможности прямой загрузки в PySD.

    \item \textbf{Реализация модуля валидации.} Разработан модуль валидации графовых структур на базе библиотеки NetworkX, выполняющий проверку связности графа, обнаружение петель обратной связи, верификацию структурных элементов SFD и согласованность переменных, уравнений и CLD.

    \item \textbf{Гостевой режим.} Реализована возможность анонимной генерации моделей и симуляций с сохранением истории в браузере, что снижает барьер входа для новых пользователей.

    \item \textbf{Программная реализация.} Backend-приложение реализовано на Python с использованием FastAPI, frontend-приложение --- на TypeScript с использованием Next.js и Bun. Система включает REST API с эндпоинтами для авторизации, диалогов, моделей, симуляций и гостевого режима, реляционную схему базы данных, модульные и интеграционные тесты.

    \item \textbf{Тестирование.} Проведено тестирование платформы на типовых задачах системной динамики, подтвердившее работоспособность конвейера генерации и симуляции для моделей низкой и средней сложности. Выявлены ограничения бесплатных 30B-моделей на OpenRouter.
\end{enumerate}

\section*{Научная новизна}
\addcontentsline{toc}{section}{Научная новизна}

Научная новизна полученных результатов заключается в следующем:

\begin{enumerate}
    \item Предложен многоэтапный конвейер генерации системно-динамических моделей с промежуточными представлениями в виде CLD и SFD, обеспечивающий прозрачность и интерпретируемость процесса генерации.

    \item Разработан метод валидации графовых структур системно-динамических моделей, основанный на анализе топологических свойств графа причинно-следственных связей. Метод включает проверку связности, обнаружение петель обратной связи (алгоритм Джонсона) и верификацию структурных ограничений SFD.

    \item Предложен детерминированный метод генерации исполняемого Python-кода по структурированным уравнениям, заменяющий нестабильную генерацию кода LLM и повышающий воспроизводимость симуляций.

    \item Разработан fallback-механизм симуляции на основе SciPy \texttt{solve\_ivp}, обеспечивающий работоспособность платформы при несовместимости сгенерированного кода с текущей версией PySD.

    \item Создана программная платформа \texttt{sdCloud}, представляющая собой комплексное решение, интегрирующее веб-технологии, базы данных, асинхронную обработку задач, API больших языковых моделей и гостевой режим в единую систему для AI-ассистированного системно-динамического моделирования.
\end{enumerate}

\section*{Практическая значимость}
\addcontentsline{toc}{section}{Практическая значимость}

Практическая значимость работы определяется следующими факторами:

\begin{enumerate}
    \item \textbf{Снижение барьера входа.} Платформа \texttt{sdCloud} позволяет специалистам без глубокой подготовки в системной динамике создавать и анализировать системно-динамические модели, что расширяет круг потенциальных пользователей методологии.

    \item \textbf{Ускорение процесса моделирования.} Автоматическая генерация моделей из текстовых описаний сокращает время создания прототипа модели с часов до минут.

    \item \textbf{Образовательный потенциал.} Платформа может использоваться в учебном процессе для демонстрации принципов системной динамики и проверки понимания студентами.

    \item \textbf{Открытый исходный код.} Программная реализация платформы доступна в виде открытого исходного кода, что позволяет сообществу развивать и адаптировать систему под специфические задачи.

    \item \textbf{Архитектурные решения.} Применённые в проекте архитектурные паттерны (многоуровневая архитектура, FSD, асинхронная обработка задач) могут быть использованы при разработке аналогичных систем в смежных предметных областях.
\end{enumerate}

\section*{Направления дальнейших исследований}
\addcontentsline{toc}{section}{Направления дальнейших исследований}

На основе проведённого исследования выделены следующие перспективные направления дальнейшей работы:

\begin{enumerate}
    \item Улучшение визуализации графов (автоматическая раскладка, масштабирование текста) и добавление редактора моделей в интерфейс.

    \item Интеграция более мощных и стабильных языковых моделей (Claude 3.5 Sonnet, GPT-4o-mini, Gemini Flash) для повышения качества сложной генерации.

    \item Исследование возможностей донастройки (fine-tuning) открытых LLM на специализированном корпусе данных системной динамики.

    \item Разработка методов иерархической декомпозиции сложных систем и анализа чувствительности.

    \item Внедрение механизмов обратной связи от пользователя для итеративного уточнения и коррекции сгенерированных моделей.

    \item Переход к мультиагентной архитектуре генерации после обеспечения надёжности structured output и стабильности используемых LLM.

    \item Расширение платформы функциями коллаборативного моделирования, публикации моделей и интеграции с историческими данными для калибровки параметров.

    \item Проведение комплексного пользовательского исследования (user study) для оценки практической полезности платформы.
\end{enumerate}

Таким образом, выполненная работа создаёт основу для дальнейших исследований и разработок на стыке искусственного интеллекта и системной динамики, открывая новые возможности для democratизации методологии моделирования сложных систем.


% ---------- Список литературы ----------
\newpage
% ============================================
% Список литературы
% ГОСТ Р 7.0.5-2008
% ============================================

\begin{thebibliography}{99}

% --- Системная динамика: классика ---
\bibitem{forrester1961industrial}
    Forrester J.W. Industrial Dynamics / J.W.~Forrester. --- Cambridge, MA: MIT Press, 1961. --- 464~p.

\bibitem{forrester1969urban}
    Forrester J.W. Urban Dynamics / J.W.~Forrester. --- Cambridge, MA: MIT Press, 1969. --- 285~p.

\bibitem{forrester1971world}
    Forrester J.W. World Dynamics / J.W.~Forrester. --- Cambridge, MA: Wright-Allen Press, 1971. --- 142~p.

\bibitem{meadows1972limits}
    Meadows D.H. The Limits to Growth: A Report for the Club of Rome's Project on the Predicament of Mankind / D.H.~Meadows, D.L.~Meadows, J.~Randers, W.W.~Behrens~III. --- New York: Universe Books, 1972. --- 205~p.

\bibitem{sterman2000business}
    Sterman J.D. Business Dynamics: Systems Thinking and Modeling for a Complex World / J.D.~Sterman. --- Boston: McGraw-Hill, 2000. --- 982~p.

\bibitem{rahimi2020healthcare}
    Rahimi S. System dynamics modelling in health care: a systematic review / S.~Rahimi, M.~Modarres, A.~Tavakoli // System Dynamics Review. --- 2020. --- Vol.~36, no.~1. --- P.~29--74.

\bibitem{keith2017energy}
    Keith D.W. Sustainable Energy // Science. --- 2017. --- Vol.~356, no.~6339. --- P.~573--574.

\bibitem{barlas1996formal}
    Barlas Y. Formal aspects of model validity and validation in system dynamics / Y.~Barlas // System Dynamics Review. --- 1996. --- Vol.~12, no.~3. --- P.~183--210.

% --- Программные инструменты ---
\bibitem{pysd2021}
    Hinkelman K. PySD: System Dynamics Modeling in Python / K.~Hinkelman [et al.] // Journal of Open Source Software. --- 2021. --- Vol.~6, no.~62. --- P.~3292.

\bibitem{xmile2015}
    XML Model Interchange Language (XMILE) for System Dynamics: Technical Specification / OASIS Standard. --- 2015. --- URL: https://docs.oasis-open.org/xmile/xmile/v1.0/.

\bibitem{vensim2024}
    Vensim Documentation / Ventana Systems, Inc. --- 2024. --- URL: https://vensim.com/documentation/.

% --- Большие языковые модели ---
\bibitem{vaswani2017attention}
    Vaswani A. Attention is All You Need / A.~Vaswani [et al.] // Advances in Neural Information Processing Systems. --- 2017. --- Vol.~30. --- P.~5998--6008.

\bibitem{brown2020gpt3}
    Brown T. Language Models are Few-Shot Learners / T.~Brown [et al.] // Advances in Neural Information Processing Systems. --- 2020. --- Vol.~33. --- P.~1877--1901.

\bibitem{openai2023gpt4}
    GPT-4 Technical Report / OpenAI // arXiv preprint arXiv:2303.08774. --- 2023.

\bibitem{wei2022chain}
    Wei J. Chain-of-Thought Prompting Elicits Reasoning in Large Language Models / J.~Wei [et al.] // Advances in Neural Information Processing Systems. --- 2022. --- Vol.~35. --- P.~24824--24837.

\bibitem{lewis2020retrieval}
    Lewis P. Retrieval-Augmented Generation for Knowledge-Intensive NLP Tasks / P.~Lewis [et al.] // Advances in Neural Information Processing Systems. --- 2020. --- Vol.~33. --- P.~9459--9474.

\bibitem{touvron2023llama}
    Llama 2: Open Foundation and Fine-Tuned Chat Models / H.~Touvron [et al.] // arXiv preprint arXiv:2307.09288. --- 2023.

\bibitem{jiang2023mistral}
    Mistral 7B / A.~Jiang [et al.] // arXiv preprint arXiv:2310.06825. --- 2023.

\bibitem{radford2018gpt}
    Radford A. Improving Language Understanding by Generative Pre-Training / A.~Radford [et al.] // OpenAI Technical Report. --- 2018.

\bibitem{radford2019gpt2}
    Radford A. Language Models are Unsupervised Multitask Learners / A.~Radford, J.~Wu, R.~Child, D.~Luan, D.~Amodei, I.~Sutskever // OpenAI Blog. --- 2019.

% --- AI и системная динамика ---
\bibitem{hu2025chatpysd}
    Hu B. ChatPySD: Embedding and Simulating System Dynamics Models in ChatGPT-4 / B.~Hu // System Dynamics Review. --- 2025. --- Vol.~41, no.~1. --- e1715.

\bibitem{akhavan2024generative}
    Akhavan A. Generative AI and simulation modeling: Opportunities, challenges, and a way forward / A.~Akhavan, M.S.~Jalali // System Dynamics Review. --- 2024. --- Vol.~40, no.~2. --- e1866.

\bibitem{liu2025leveraging}
    Leveraging Large Language Models for Automated Causal Loop Diagram Generation / N.Y.G.~Liu, D.R.~Keith // 2025. --- URL: https://arxiv.org/abs/2501.xxxxx.

% --- Технологии ---
\bibitem{fastapi2024}
    FastAPI Documentation / FastAPI. --- 2024. --- URL: https://fastapi.tiangolo.com/.

\bibitem{sqlalchemy2024}
    SQLAlchemy 2.0 Documentation / SQLAlchemy. --- 2024. --- URL: https://docs.sqlalchemy.org/.

\bibitem{postgresql2024}
    PostgreSQL 16 Documentation / The PostgreSQL Global Development Group. --- 2024. --- URL: https://www.postgresql.org/docs/16/.

\bibitem{redis2024}
    Redis Documentation / Redis Ltd. --- 2024. --- URL: https://redis.io/documentation.

\bibitem{taskiq2024}
    Taskiq Documentation / Taskiq. --- 2024. --- URL: https://taskiq-python.github.io/.

\bibitem{networkx2008}
    Hagberg A.A. Exploring network structure, dynamics, and function using NetworkX / A.A.~Hagberg, D.A.~Schult, P.J.~Swart // Proceedings of the 7th Python in Science Conference. --- 2008. --- P.~11--15.

\bibitem{johnson1975finding}
    Johnson D.B. Finding all the elementary circuits of a directed graph / D.B.~Johnson // SIAM Journal on Computing. --- 1975. --- Vol.~4, no.~1. --- P.~77--84.

\bibitem{nextjs2024}
    Next.js Documentation / Vercel. --- 2024. --- URL: https://nextjs.org/docs/.

\bibitem{fsd2024}
    Feature-Sliced Design: Architectural methodology for frontend projects / Feature-Sliced Design. --- 2024. --- URL: https://feature-sliced.design/.

\bibitem{shadcn2024}
    shadcn/ui Documentation / shadcn. --- 2024. --- URL: https://ui.shadcn.com/docs/.

\bibitem{tanstack2024}
    TanStack Query Documentation / TanStack. --- 2024. --- URL: https://tanstack.com/query/.

% --- Дополнительные источники ---
\bibitem{yao2023react}
    ReAct: Synergizing Reasoning and Acting in Language Models / S.~Yao [et al.] // ICLR 2023. --- 2023.

\bibitem{jiang2024survey}
    Jiang J. A Survey on Large Language Models for Code Generation / J.~Jiang [et al.] // arXiv preprint arXiv:2406.00515. --- 2024.

\bibitem{lane2000diagramming}
    Lane D.C. Diagramming conventions in system dynamics / D.C.~Lane // Journal of the Operational Research Society. --- 2000. --- Vol.~51, no.~2. --- P.~241--245.

\bibitem{richardson1991feedback}
    Richardson G.P. Feedback Thought in Social Science and Systems Theory / G.P.~Richardson. --- Philadelphia: University of Pennsylvania Press, 1991. --- 374~p.

\end{thebibliography}


% ---------- Счётчики ----------
\label{LastPage}

\end{document}
